\chapter{Application to Theorems of Good Reduction}
\label{ch:14}

\section{General review of vanishing cycles (cf.\ \cite[7, Exp.~XIII, 2]{SGA})}
\label{sec:14.1}

\kmitem\label{item:14.1.1}
Fix a prime number $p$.
Let $V$ be a strictly henselian discrete valuation ring with algebraically closed residue field $k$ of characteristic $p$, and fraction field $K$. Let $\overline{K}$ be a separable closure of $K$. We follow the standard geometric notations
\begin{equation}\label{eq:14.1.1.1}
s = \Spec(k) \mathrel{\lhook\joinrel\longrightarrow} S = \Spec(V) \hookleftarrow \bar{\eta} = \Spec(\overline{K}) .
\end{equation}
Fix a finite extension $E/\mathbb{Q}$, and a finite place $\lambda$ of $E$ of residue characteristic $\ell \neq p$. Let
\begin{equation}\label{eq:14.1.1.2}
\begin{tikzcd}[row sep=1.8em]
Z \arrow[d, "f"] \\
S = \Spec(V)
\end{tikzcd}
\end{equation}
be a proper morphism, and let $\mathcal{F}$ be a constructible $E_\lambda$-sheaf on $Z$. Form the cartesian diagram
\begin{equation}\label{eq:14.1.1.3}
\begin{tikzcd}[column sep=2.5em, row sep=2.2em]
\mathcal{F}_s \overset{\mathrm{dfn}}{=} i^*(\mathcal{F})\ ;\ Z_s \arrow[r, "i"] \arrow[d] & Z \arrow[d] & Z_{\bar{\eta}}\ ;\ \mathcal{F}_{\bar{\eta}} \overset{\mathrm{dfn}}{=} (\bar{j})^*(\mathcal{F}) \arrow[l, hook', "\bar{j}"'] \arrow[d] \\
s \arrow[r, hook] & S & \bar{\eta} \arrow[l, hook'] .
\end{tikzcd}
\end{equation}

\kmitem\label{item:14.1.2}
The derived category version of the Leray spectral sequence for $\bar{j}$ may be written
\begin{equation}\label{eq:14.1.2.1}
\begin{tikzcd}[row sep=1.8em]
R\Gamma(Z_{\bar{\eta}}, \mathcal{F}_{\bar{\eta}}) = R\Gamma(Z, R(\bar{j})_* \mathcal{F}_{\bar{\eta}}) \arrow[d, "\rotatebox{90}{$\sim$}"] \\
R\Gamma(Z_s, i^* R(\bar{j})_* \mathcal{F}_{\bar{\eta}}) ,
\end{tikzcd}
\end{equation}
the vertical isomorphism because $Z$ is proper over $S$. Because $\mathcal{F}_{\bar{\eta}}$ is by definition $(\bar{j})^*(\mathcal{F})$, the natural morphism of adjunction yields an arrow
\begin{equation}\label{eq:14.1.2.2}
\begin{tikzcd}[row sep=2.2em, column sep=3em]
\mathcal{F} \arrow[r] \arrow[dr, dashed] & \bar{j}_* (\bar{j})^* \mathcal{F} = \bar{j}_* \mathcal{F}_{\bar{\eta}} \arrow[d] \\
& R(\bar{j})_* \mathcal{F}_{\bar{\eta}} ,
\end{tikzcd}
\end{equation}
whence, by restriction to the special fiber, a morphism of complexes on $Z_s$
\begin{equation}\label{eq:14.1.2.3}
\mathcal{F}_s = i^* \mathcal{F} \to i^* R(\bar{j})_* \mathcal{F}_{\bar{\eta}} .
\end{equation}
By definition, its mapping cone, with $\mathcal{F}_s$ viewed as lying in degree $-1$
\begin{equation}\label{eq:14.1.2.4}
\mathcal{F}_s \to i^* R(\bar{j})_* \mathcal{F}_{\bar{\eta}}
\end{equation}
is denoted
\begin{equation}\label{eq:14.1.2.5}
R\Phi(\mathcal{F}), \quad\text{or}\quad R\Phi_{Z/S}(\mathcal{F}) .
\end{equation}
The complex $R\Psi_{Z/S}(\mathcal{F})$ on $Z_s$ is defined by
\begin{equation}\label{eq:14.1.2.6}
R\Psi_{Z/S}(\mathcal{F}) \overset{\mathrm{dfn}}{=} i^* R(\bar{j})_* \mathcal{F}_{\bar{\eta}} .
\end{equation}
The short exact sequence of complexes on $Z_s$
\begin{equation}\label{eq:14.1.2.7}
0 \to R\Psi_{Z/S}(\mathcal{F}) \to R\Phi_{Z/S}(\mathcal{F}) \to (\mathcal{F}_s \text{, placed in degree } -1) \to 0
\end{equation}
gives an exact sequence of sheaves on $Z_s$
\begin{equation}\label{eq:14.1.2.8}
0 \to R^{-1}\Phi_{Z/S}(\mathcal{F}) \to \mathcal{F}_s \to i^*(\bar{j})_* \mathcal{F}_{\bar{\eta}} \to R^0 \Phi_{Z/S}(\mathcal{F}) \to 0
\end{equation}
and, for all integers $i \ge 1$, an isomorphism of sheaves on $Z_s$
\begin{equation}\label{eq:14.1.2.9}
R^i \Psi_{Z/S}(\mathcal{F}) \xrightarrow{\ \sim\ } R^i \Phi_{Z/S}(\mathcal{F}) .
\end{equation}
Passing to cohomology on $Z_s$, the above exact sequence of complexes gives rise to a long exact cohomology sequence
\begin{equation}\label{eq:14.1.2.10}
\begin{tikzcd}[column sep=0.55em, font=\small]
\dots \arrow[r] & H^i(Z_s, \mathcal{F}_s) \arrow[r] & H^i(Z_s, R\Psi_{Z/S}(\mathcal{F})) \arrow[r] \arrow[d, equal] & H^i(Z_s, R\Phi_{Z/S}(\mathcal{F})) \arrow[r] & \dots \\
& & H^i(Z_{\bar{\eta}}, \mathcal{F}_{\bar{\eta}}) & &
\end{tikzcd}
\end{equation}
Thus it is only the cohomology groups
\begin{equation}\label{eq:14.1.2.11}
H^i(Z_s, R\Phi_{Z/S}(\mathcal{F}))
\end{equation}
which prevent the specialization map
\begin{equation}\label{eq:14.1.2.12}
H^i(Z_s, \mathcal{F}_s) \to H^i(Z_{\bar{\eta}}, \mathcal{F}_{\bar{\eta}})
\end{equation}
from being an isomorphism for all $i$.

\kmitem\label{item:14.1.3}\emph{The inertia group}
\[
I_K = \operatorname{Gal}(\overline{K}/K)
\]
operates naturally on the entire situation, trivially on $(Z, \mathcal{F})$ and on $(Z_s, \mathcal{F}_s)$, highly non-trivially on $(Z_{\bar{\eta}}, \mathcal{F}_{\bar{\eta}})$. In particular the above exact sequence is $I_K$-equivariant, with $I_K$ operating trivially on $H^i(Z_s, \mathcal{F}_s)$.

\begin{definition}\label{def:14.1.4}
Let $L/K$ be a finite separable extension of $K$. An $E_\lambda$-subspace of $H^i(Z_{\bar{\eta}}, \mathcal{F}_{\bar{\eta}})$ is said to have ``good reduction over $L$'' if it is stable by the subgroup $I_L \subset I_K$, and if $I_L$ operates trivially on it.
\par
For example, any subspace of
\[
\Image(H^i(Z_s, \mathcal{F}_s) \to H^i(Z_{\bar{\eta}}, \mathcal{F}_{\bar{\eta}}))
\]
has good reduction over $K$ itself. We will be particularly interested in the case $i = 1$. So it will be very useful to know when the specialization map
\[
H^1(Z_s, \mathcal{F}_s) \to H^1(Z_{\bar{\eta}}, \mathcal{F}_{\bar{\eta}})
\]
is injective.
\end{definition}

\begin{klem}\label{lem:14.1.5}
Suppose we have a commutative diagram
\[
\begin{tikzcd}[row sep=1.8em, column sep=1.5em]
\mathcal{F}\ ;\ Z \arrow[dr, "f"'] & & W\ ;\ \pi^*\mathcal{F} \arrow[ll, "\pi"'] \arrow[dl, "g"] \\
& S &
\end{tikzcd}
\]
with $(\mathcal{F}, Z, f, S)$ as in (\ref{item:14.1.1}) above, and $\pi$ finite. If $\pi_{\bar{\eta}} : W_{\bar{\eta}} \to Z_{\bar{\eta}}$ is an isomorphism, then we have canonically on $Z_s$,
\[
R\Psi_{Z/S}(\mathcal{F}) \simeq (\pi_s)_* R\Psi_{W/S}(\pi^*\mathcal{F}) .
\]
\end{klem}

\begin{proof}
We have a commutative diagram with both squares cartesian
\begin{center}
\begin{tikzcd}[row sep=2em, column sep=2.5em]
W_s \arrow[r, hook, "i_W"] \arrow[d, "\pi_s"'] & W \arrow[d, "\pi"] & W_{\bar{\eta}} \arrow[l, hook', "\bar{j}_W"'] \arrow[d, "\pi_{\bar{\eta}}"', "\rotatebox{90}{$\sim$}"] \\
Z_s \arrow[r, hook, "i_Z"'] & Z & Z_{\bar{\eta}} \arrow[l, hook', "\bar{j}_Z"] .
\end{tikzcd}
\end{center}
By definition,
\[
(\pi_s)_* R\Psi_{W/S}(\pi^*\mathcal{F}) = (\pi_s)_* (i_W)^* R(\bar{j}_W)_* ((\bar{j}_W)^* \pi^*\mathcal{F})
\]
which, by the proper base-change theorem applied to $(\pi_s)_* = R(\pi_s)_*$, gives
\[
\begin{aligned}
&= (i_Z)^* R(\pi)_* R(\bar{j}_W)_* ((\pi_{\bar{\eta}})^* (\bar{j}_Z)^* \mathcal{F}) \\
&= (i_Z)^* R(\bar{j}_Z)_* R(\pi_{\bar{\eta}})_* (\pi_{\bar{\eta}})^* (\bar{j}_Z)^* \mathcal{F} \\
&= (i_Z)^* R(\bar{j}_Z)_* ((\bar{j}_Z)^* \mathcal{F}) = R\Psi_{Z/S}(\mathcal{F}) . \qedhere
\end{aligned}
\]
\end{proof}

\begin{klem}\label{lem:14.1.6}
In the situation of the previous lemma (\ref{lem:14.1.5}), suppose that a finite group $G$ acts $k$-linearly on the special fiber $W_s$ of $W$, that $G$ acts trivially on $Z_s$, and that the morphism $\pi_s : W_s \to Z_s$ is $G$-equivariant, and induces by passage to quotients a radicial surjective integral morphism
\[
W_s/G \to Z_s .
\]
Then we have canonical isomorphisms on $Z_s$
\begin{enumerate}[label=\arabic*)]
\item $\mathcal{F}_s \xrightarrow{\ \sim\ } ((\pi_s)_* (\pi_s)^* \mathcal{F}_s)^G$
\item $R\Psi_{Z/S}(\mathcal{F}) \xrightarrow{\ \sim\ } ((\pi_s)_* R\Psi_{W/S}(\pi^*\mathcal{F}))^G$
\item $R\Phi_{Z/S}(\mathcal{F}) \xrightarrow{\ \sim\ } ((\pi_s)_* R\Phi_{W/S}(\pi^*\mathcal{F}))^G .$
\end{enumerate}
\end{klem}

\begin{proof}
The first assertion holds because the map $W_s/G \to Z_s$ is radicial, surjective and integral. By the definition of $R\Phi$ as a mapping cone, assertion 3) follows easily once we know 1) and 2). To prove 2), we compute, using the diagram
\begin{center}
\begin{tikzcd}[row sep=2em, column sep=2.5em]
W_s \arrow[r, hook, "i_W"] \arrow[d, "\pi_s"'] & W \arrow[d, "\pi"] & W_{\bar{\eta}} \arrow[l, hook', "\bar{j}_W"'] \arrow[d, "\pi_{\bar{\eta}}"] \\
Z_s \arrow[r, hook, "i_Z"'] & Z & Z_{\bar{\eta}} \arrow[l, hook', "\bar{j}_Y"]
\end{tikzcd}
\end{center}
\begin{align*}
(\pi_s)_* R\Psi_{X/S}(\pi^*\mathcal{F}) &= (\pi_s)_* (i_X)^* R(\bar{j}_W)_* ((\bar{j}_W)^*\pi^*\mathcal{F}) \\
&= (i_Z)^* \pi_* R(\bar{j}_W)_* ((\pi_{\bar{\eta}})^* (\bar{j}_Z)^*(\mathcal{F})) \\
&= (i_Z)^* R(\bar{j}_Z)_* (\pi_{\bar{\eta}})_* (\pi_{\bar{\eta}})^* (\bar{j}_Z)^*(\mathcal{F}) .
\end{align*}
Applying the exact functor ``$G$-invariants'', we find
\[
((\pi_s)_* R\Psi_{W/S}(\pi^*\mathcal{F}))^G = (i_Z)^* R(\bar{j}_Z)_* (((\pi_{\bar{\eta}})_* (\pi_{\bar{\eta}})^* (\bar{j}_Z)^*(\mathcal{F}))^G) .
\]
Because $\pi_{\bar{\eta}} : W_{\bar{\eta}} \to Z_{\bar{\eta}}$ makes $W_{\bar{\eta}}/G \xrightarrow{\sim} Z_{\bar{\eta}}$, we have, for any constructible $E_\lambda$-sheaf $\mathcal{G}$ on $Z_{\bar{\eta}}$,
\[
\mathcal{G} \xrightarrow{\ \sim\ } ((\pi_{\bar{\eta}})_* (\pi_{\bar{\eta}})^* \mathcal{G})^G .
\]
Taking for $\mathcal{G}$ the sheaf $(\bar{j}_Z)^*(\mathcal{F})$, we find the required isomorphism 2).
\end{proof}

\begin{kcor}\label{cor:14.1.7}
In the situation of the previous lemma (\ref{lem:14.1.6}), the two specialization maps
\begin{align*}
H^i(Z_s, \mathcal{F}_s) &\to H^i(Z_{\bar{\eta}}, \mathcal{F}_{\bar{\eta}}) \\
H^i(W_s, (\pi_s)^*\mathcal{F}_s) &\to H^i(W_{\bar{\eta}}, \mathcal{F}_{\bar{\eta}})
\end{align*}
sit in a commutative diagram
\begin{center}
\begin{tikzcd}[column sep=3em, row sep=2em]
H^i(Z_s, \mathcal{F}_s) \arrow[r] \arrow[d, equal] & H^i(Z_{\bar{\eta}}, \mathcal{F}_{\bar{\eta}}) \\
H^i(Z_s, ((\pi_s)_* (\pi_s)^*\mathcal{F}_s)^G) \arrow[d, equal] & \\
H^i(W_s, (\pi_s)^*\mathcal{F}_s)^G \arrow[d, hook] & \\
H^i(W_s, (\pi_s)^*\mathcal{F}_s) \arrow[r] & H^i(W_{\bar{\eta}}, \mathcal{F}_{\bar{\eta}}) \arrow[uuu, "(\pi_{\bar{\eta}})^*"', "\rotatebox{90}{$\sim$}"] .
\end{tikzcd}
\end{center}
\end{kcor}

\begin{kcor}\label{cor:14.1.8}
In the situation of the previous lemma (\ref{lem:14.1.6}), if the specialization map
\[
H^i(W_s, (\pi_s)^*\mathcal{F}_s) \to H^i(W_{\bar{\eta}}, \mathcal{F}_{\bar{\eta}})
\]
is injective, then the specialization map
\[
H^i(Z_s, \mathcal{F}_s) \to H^i(Z_{\bar{\eta}}, \mathcal{F}_{\bar{\eta}})
\]
is also injective.
\end{kcor}

\kmitem\label{item:14.1.9}
We now consider the effect of extension of the ground field. Thus we fix a finite extension $L/K$, and denote by $V_L$ the integral closure of $V$ in $L$. Thus $V_L$ is itself a strictly henselian discrete valuation ring, finite over $V$, with the same residue field $k$ as $V$ itself. Let us write
\[
t = \Spec(k) \mathrel{\lhook\joinrel\longrightarrow} T = \Spec(V_L) \hookleftarrow \bar{\eta} = \Spec(\overline{K}) ,
\]
and form the cartesian diagram
\begin{center}
\begin{tikzcd}[row sep=2em, column sep=2.5em]
\mathcal{F}\ ;\ Z \arrow[d, "f"'] & Z_T\ ;\ \lambda^*\mathcal{F} \arrow[l, "\lambda"'] \arrow[d] \\
S & T \arrow[l] .
\end{tikzcd}
\end{center}

\begin{klem}\label{lem:14.1.9.1}
In the above situation (\ref{item:14.1.9}), $Z_s = (Z_T)_t$, and we have natural equalities
\begin{enumerate}[label=\arabic*)]
\item $(\lambda^*\mathcal{F})_t = \lambda^*\mathcal{F}|(Z_T)_t = \mathcal{F}|Z_s = \mathcal{F}_s .$
\item $R\Psi_{Z_T/T}(\lambda^*\mathcal{F}) = R\Psi_{Z/S}(\mathcal{F}) .$
\item $R\Phi_{Z_T/T}(\lambda^*\mathcal{F}) = R\Phi_{Z/S}(\mathcal{F}) .$
\end{enumerate}
\end{klem}

\begin{proof}
Consider $\lambda$ as an $S$-morphism
\begin{center}
\begin{tikzcd}[row sep=1.8em, column sep=1.5em]
Z & & Z_T \arrow[ll, "\lambda"'] \arrow[dl, dashed] \\
& S &
\end{tikzcd}
\end{center}
The map $\lambda_s : (Z_T)_s \to Z_s$ is surjective, radicial and finite, because $V_L/\mathfrak{m}_K V_L$ is artin local with residue field $k$. But $S$ and $T$ have the same $\bar{\eta} = \Spec(\overline{K})$, so we have a commutative diagram
\begin{center}
\begin{tikzcd}[row sep=2.2em, column sep=2.5em]
& & Z_{\bar{\eta}} \arrow[dll, "\bar{j}_Z"', bend right=20] \arrow[dl, "\bar{j}_{Z_T}"] \\
Z & Z_T \arrow[l, "\lambda"'] & \\
Z_s = (Z_T)_t \arrow[u, "i_Z"] & (Z_T)_s \arrow[l, "\lambda_s"] \arrow[u, "i_{Z_T}"'] &
\end{tikzcd}
\end{center}
Thus we readily calculate
\begin{align*}
R\Psi_{Z/S}(\mathcal{F}) &= (i_Z)^* R(\bar{j}_Z)_* ((\bar{j}_Z)^*\mathcal{F}) \\
&= (i_Z)^* R\lambda_* R(\bar{j}_{Z_T})_* ((\bar{j}_{Z_T})^*(\lambda^*\mathcal{F})) \\
&= R(\lambda_s)_* (i_{Z_T})^* R(\bar{j}_{Z_T})_* ((\bar{j}_{Z_T})^*(\lambda^*\mathcal{F})) \\
&= R(\lambda_s)_* (\lambda_s)^* (R\Psi_{Z_T/T}(\lambda^*\mathcal{F})) \\
&= R\Psi_{Z_T/T}(\lambda^*\mathcal{F}) .
\end{align*}
The third from last equality is proper base change for $\lambda$. The last two equalities hold because $\lambda_s$ is surjective, radicial and finite, so $R(\lambda_s)_* = (\lambda_s)_*$ and $\lambda_s^*$ are inverse equivalences of categories.
\par
Similarly,
\[
\begin{tikzcd}[row sep=1.8em]
R\Phi_{Z/S}(\mathcal{F}) = [\mathcal{F}|Z_s \to R\Psi_{Z/S}(\mathcal{F})] \arrow[d, equal] \\
R\Phi_{Z_T/T}(\mathcal{F}) = [\lambda^*\mathcal{F}|(Z_T)_t \to R\Psi_{Z_T/T}(\lambda^*\mathcal{F})] .
\end{tikzcd}
\qedhere
\]
\end{proof}

\kmitem\label{item:14.1.10}
We now consider the effect of a finite group of operators.

\begin{klem}\label{lem:14.1.10.1}
Let $V$ be a strictly henselian discrete valuation ring, $X$ a projective scheme over $S = \Spec(V)$ on which a finite group $G$ operates $V$-linearly, $Y$ the quotient $X/G$, and $\mathcal{F}$ a constructible $E_\lambda$-sheaf on $Y$:
\[
\begin{tikzcd}[row sep=1.8em, column sep=1.5em]
\mathcal{F}\ ;\ Y = X/G \arrow[dr] & & X\ ;\ \pi^*\mathcal{F} \arrow[ll, "\pi"'] \arrow[dl] \\
& S . &
\end{tikzcd}
\]
Then we have canonical isomorphisms on $Y_s$
\begin{enumerate}[label=\arabic*)]
\item $\mathcal{F}_s = ((\pi_s)_*(\pi_s)^*\mathcal{F}_s)^G$
\item $R\Psi_{Y/S}(\mathcal{F}) = ((\pi_s)_* R\Psi_{X/S}(\pi^*\mathcal{F}))^G$
\item $R\Phi_{Y/S}(\mathcal{F}) = ((\pi_s)_* R\Phi_{X/S}(\pi^*\mathcal{F}))^G .$
\end{enumerate}
\end{klem}

\begin{proof}
The first assertion holds because the map $X_s/G \to Y_s$ is radicial, surjective and integral. By the definition of $R\Phi$ as a mapping cone, assertion 3) follows easily once we know 1) and 2). To prove 2), we compute, using the diagram
\begin{center}
\begin{tikzcd}[row sep=2em, column sep=2.5em]
X_s \arrow[r, hook, "i_X"] \arrow[d, "\pi_s"'] & X \arrow[d, "\pi"] & X_{\bar{\eta}} \arrow[l, hook', "\bar{j}_X"'] \arrow[d, "\pi_{\bar{\eta}}"] \\
Y_s \arrow[r, hook, "i_Y"'] & Y & Y_{\bar{\eta}} \arrow[l, hook', "\bar{j}_Y"]
\end{tikzcd}
\end{center}
\begin{align*}
(\pi_s)_* R\Psi_{X/S}(\pi^*\mathcal{F}) &= (\pi_s)_* (i_X)^* R(\bar{j}_X)_* ((\bar{j}_X)^*\pi^*\mathcal{F}) \\
&= (i_Y)^* \pi_* R(\bar{j}_X)_* ((\pi_{\bar{\eta}})^* (\bar{j}_Y)^*(\mathcal{F})) \\
&= (i_Y)^* R(\bar{j}_Y)_* (\pi_{\bar{\eta}})_* (\pi_{\bar{\eta}})^* (\bar{j}_Y)^*(\mathcal{F}) .
\end{align*}
Applying the exact functor ``$G$-invariants'', we find
\[
((\pi_s)_* R\Psi_{X/S}(\pi^*\mathcal{F}))^G = (i_Y)^* R(\bar{j}_Y)_* (((\pi_{\bar{\eta}})_* (\pi_{\bar{\eta}})^* (\bar{j}_Y)^*(\mathcal{F}))^G) .
\]
Because $\pi_{\bar{\eta}} : X_{\bar{\eta}} \to Y_{\bar{\eta}}$ makes $X_{\bar{\eta}}/G \xrightarrow{\sim} Y_{\bar{\eta}}$, we have, for any constructible $E_\lambda$-sheaf $\mathcal{G}$ on $Y_{\bar{\eta}}$,
\[
\mathcal{G} \xrightarrow{\ \sim\ } ((\pi_{\bar{\eta}})_* (\pi_{\bar{\eta}})^* \mathcal{G})^G .
\]
Taking for $\mathcal{G}$ the sheaf $(\bar{j}_Y)^*(\mathcal{F})$, we find the required isomorphism 2).
\end{proof}

\begin{kcor}\label{cor:14.1.11}
Hypotheses as in (\ref{lem:14.1.10.1}) above, the long exact cohomology sequence with its $I_K$-action
\begin{multline*}
\dots \to H^i(Y_s, \mathcal{F}_s) \to H^i(Y_{\bar{\eta}}, \mathcal{F}_{\bar{\eta}}) \\
\to H^i(Y_s, R\Phi_{Y/S}(\mathcal{F})) \to H^{i+1}(Y_s, \mathcal{F}_s) \to \dots
\end{multline*}
is canonically isomorphic to the $G$-invariants in the long exact cohomology sequence
\begin{multline*}
\dots \to H^i(X_s, (\pi_s)^*\mathcal{F}_s) \to H^i(X_{\bar{\eta}}, (\pi_{\bar{\eta}})^*\mathcal{F}_{\bar{\eta}}) \\
\to H^i(X_s, R\Phi_{X/S}(\pi^*\mathcal{F})) \to H^{i+1}(X_s, (\pi_s)^*\mathcal{F}_s) \to \dots .
\end{multline*}
\end{kcor}

\begin{kcor}\label{cor:14.1.12}
Hypotheses as in (\ref{lem:14.1.10.1}) above, suppose that
\[
R^i \Phi_{X/S}(\pi^*\mathcal{F}) = \begin{cases}
0 & \text{for } i \neq 1 \\
\text{a punctual sheaf,} & \text{for } i = 1 .
\end{cases}
\]
Then we have
\[
R^i \Phi_{Y/S}(\mathcal{F}) = \begin{cases}
0 & \text{for } i \neq 1 \\
\text{a punctual sheaf,} & \text{for } i = 1 .
\end{cases}
\]
In particular, the specialization maps
\[
H^1(Y_s, \mathcal{F}_s) \to H^1(Y_{\bar{\eta}}, \mathcal{F}_{\bar{\eta}})
\]
and
\[
H^1(X_s, \pi_s^*\mathcal{F}_s) \to H^1(X_{\bar{\eta}}, (\pi_{\bar{\eta}})^*\mathcal{F}_{\bar{\eta}})
\]
are both injective.
\end{kcor}

\section{Application to curves}
\label{sec:14.2}

\kmitem\label{item:14.2.1}
We now suppose that $V$ is a mixed characteristic strictly henselian discrete valuation ring with algebraically closed residue field $k$ of characteristic $p$, and that we are given the following geometric data once and for all.
\begin{enumerate}[label=\arabic*)]
\item a proper and smooth curve with geometrically connected fibers
\[
\begin{tikzcd}[row sep=1.8em]
Y \arrow[d] \\
S = \Spec(V) .
\end{tikzcd}
\]
\item a closed subscheme (``the cusps'') $D \subset Y$, finite etale over $S$.
\item a finite extension $E/\mathbb{Q}$, a finite place $\lambda$ of $E$ of residue characteristic $\ell \neq p$, and a constructible $E_\lambda$-sheaf $\mathcal{F}$ on $Y$, which is \emph{lisse} when restricted to $D$.
\end{enumerate}
Notice that, $V$ being of generic characteristic zero, $\mathcal{F}$ is automatically \emph{tamely ramified} along $D$.

\kmitem\label{item:14.2.2}
We will try to make do with the following convention; given any $Y$-scheme
\[
\begin{tikzcd}[row sep=1.8em]
X \arrow[d, "g"] \\
Y ,
\end{tikzcd}
\]
we will denote simply by $\mathcal{F}$ the sheaf $g^*(\mathcal{F})$ on $X$, (the usual convention when dealing with the constant sheaf).

\kmitem\label{item:14.2.3}
Now suppose we are given
\begin{enumerate}[label=\arabic*), start=4]
\item a finite galois extension $L/K$, with galois group $G$;
\item a finite $Y \tens{V} V_L$-scheme which is normal
\[
\begin{tikzcd}[row sep=1.8em]
X \arrow[d, "\rho"] \\
Y \tens{V} V_L ,
\end{tikzcd}
\]
\item a finite group $\Gamma$, given with a surjective homomorphism
\[
\Gamma \xrightarrow{\text{``}\det\text{''}} G ,
\]
and a $Y$-linear action of $\Gamma$ on $X$ such that $\gamma \in \Gamma$ acts $\det(\gamma)$-semilinearly. We denote by $\mathrm{SL} \subset \Gamma$ the subgroup of elements having trivial ``$\det$'': $\Gamma \twoheadrightarrow G$.
\end{enumerate}

\begin{theorem}\label{thm:14.2.4}
Hypotheses as in \ref{item:14.2.1}--\ref{item:14.2.3} above, suppose that $X$ satisfies the following conditions:
\begin{enumerate}[label=\arabic*)]
\item $X$ is flat over $V_L$
\item outside a finite set of $k$-valued ``supersingular points'' in the special fiber, all of which lie over $Y - D$, $X$ is a smooth curve over $V_L$
\item the special fiber $X \tens{V_L} k$ is reduced
\item the scheme of cusps of $X$, defined as $(\rho^{-1}(D))^{\mathrm{red}}$, is finite etale over $V_L$.
\end{enumerate}
Then we have
\begin{enumerate}[label=\arabic*)]
\item The sheaves $R^i \Phi_{X/V_L}(\mathcal{F})$ vanish for $i \neq 1$, and $R^1 \Phi_{X/V_L}(\mathcal{F})$ is a punctual sheaf, supported at the supersingular points.
\item For any subgroup $\Gamma_1 \subset \Gamma$, denote by $V_1 \subset V_L$ the subring $(V_L)^{\det(\Gamma_1)}$, and by $X_1$ the quotient $X/\Gamma_1$. Then $X_1$ is a normal scheme, finite over $Y \tens{V} V_1$, and the specialization map
\[
H^1((X/\Gamma_1)_s, \mathcal{F}) \to H^1((X/\Gamma_1)_{\bar{\eta}}, \mathcal{F})
\]
is injective.
\item For any subgroup $\Gamma_1 \subset \Gamma$, the specialization maps for $X/\Gamma_1$ over $V_1$ and for $X$ over $V_L$ are related by an $I_L$-equivariant commutative diagram
\[
\begin{tikzcd}[row sep=1.6em, column sep=2.5em]
H^1((X/\Gamma_1)_s, \mathcal{F}) \arrow[r, hook] \arrow[d, phantom, "\wr"] & H^1((X/\Gamma_1)_{\bar{\eta}}, \mathcal{F}) \arrow[d, phantom, "\wr"] \\
H^1(X_s/\Gamma_1, \mathcal{F}) \arrow[d, phantom, "\wr"] & H^1(X_{\bar{\eta}}/\mathrm{SL} \cap \Gamma_1, \mathcal{F}) \arrow[d, phantom, "\wr"] \\
H^1(X_s, \mathcal{F})^{\Gamma_1} \arrow[d, hook] & H^1(X_{\bar{\eta}}, \mathcal{F})^{\mathrm{SL} \cap \Gamma_1} \arrow[d, hook] \\
H^1(X_s, \mathcal{F}) \arrow[r, hook] & H^1(X_{\bar{\eta}}, \mathcal{F}) .
\end{tikzcd}
\]
\item If $\Gamma_2$ is any subgroup of $\Gamma$ which normalizes $\Gamma_1$ and which has ``$\det$''$(\Gamma_2)$ trivial, then $\Gamma_2$ operates naturally and $I_{L^{\det(\Gamma_1)}}$-equivariantly (resp.\ $I_L$-equivariantly) on the top (resp.\ bottom) row of the above diagram, and the entire diagram is $\Gamma_2$-equivariant.
\end{enumerate}
\end{theorem}

\begin{proof}
We begin with 1). The sheaf $\mathcal{F}$ on $X$ is lisse on $X - \{\text{cusps}\}$, and it is lisse when restricted to the cusps. Because we are over a base of generic characteristic zero, and the cusps are finite etale over the base, $\mathcal{F}$ is automatically tamely ramified along the cusps of $X$. Therefore (\cite[7, Exp.~XIII, 2.1.11]{SGA}) all the $R^i \Phi_{X/V_L}(\mathcal{F})$ vanish along the cusps.
Over $X - \{\text{cusps}\}$, $\mathcal{F}$ is lisse, and $X - \{\text{cusps}\}$ is lisse over the base $\Spec(V_L)$ outside a finite set of ``supersingular'' points in the special fiber. Therefore all the sheaves $R^i \Phi_{X/V_L}(\mathcal{F})$ are punctual, supported at the supersingular points. For dimension reasons, $R^i \Phi$ vanishes for $i \ge 2$. To see that $R^{-1}\Phi$ and $R^0 \Phi$ vanish at the supersingular points, it suffices, $\mathcal{F}$ being lisse there, to show that $X \to \Spec(V_L)$ is locally zero-acyclic (\cite[4, Exp.~XV, 1.11]{SGA}) but this follows from its being flat with reduced geometric fibers (\cite[4, Exp.~XV, 4.1]{SGA}). This proves assertion 1).

We prove assertions 2) and 3) together. Consider the diagram
\[
\begin{tikzcd}[row sep=3.5em, column sep=1.2em]
X/\Gamma_1 \arrow[d] & (X/\Gamma_1) \otimes V_L \arrow[l] \arrow[d] & X/(\mathrm{SL} \cap \Gamma_1) \arrow[l] \arrow[ld] & X \arrow[l] \arrow[lld] \\
\Spec((V_L)^{\det(\Gamma_1)}) & \Spec(V_L) \arrow[l] & & {}
\arrow[phantom, from=1-1, to=2-2, "\text{\textcircled{\scriptsize 1}}"]
\arrow[phantom, from=1-2, to=1-3, shift right=2em, "\text{\textcircled{\scriptsize 2}}"]
\arrow[phantom, from=1-3, to=1-4, shift right=2em, "\text{\textcircled{\scriptsize 3}}"]
\end{tikzcd}
\]
We apply (\ref{cor:14.1.11}, \ref{cor:14.1.12}) to the triangle marked \text{\textcircled{\scriptsize 3}}, then we apply (\ref{cor:14.1.7}) to the triangle marked \text{\textcircled{\scriptsize 2}}, the group in this case being $\Gamma_1/\mathrm{SL} \cap \Gamma_1 \xrightarrow{\ \sim\ } \det(\Gamma_1)$, and finally we apply (\ref{lem:14.1.9.1}) to the cartesian square marked \text{\textcircled{\scriptsize 1}}.

Assertion 4) is clear by functoriality.
\end{proof}

\begin{corollary}[Numerical criterion for good reduction]\label{cor:14.2.5}
Let $\Gamma_2$ be a subgroup of $\Gamma$ which normalizes $\Gamma_1$ and has $\det(\Gamma_2)$ trivial. Let $\Sigma$ be any $E_\lambda$-rational set of $\overline{E}_\lambda$-representations of $\Gamma_2$. For any finite-dimensional $E_\lambda$-representation $M$ of $\Gamma_2$, let us denote by $M(\Sigma)$ the $E_\lambda$-subrepresentation of $M$ consisting of the sum of all isotypical components of type $\sigma \in \Sigma$. Then
\begin{enumerate}[label=\arabic*)]
\item We always have an inequality of dimensions
\[
\dim H^1(X_s, \mathcal{F})^{\Gamma_1}(\Sigma) \le \dim H^1(X_{\bar{\eta}}, \mathcal{F})^{\mathrm{SL} \cap \Gamma_1}(\Sigma) .
\]
\item If this inequality is an equality, then the inclusion
\[
\begin{tikzcd}[row sep=1.8em, column sep=1.8em]
H^1(X_s, \mathcal{F})^{\Gamma_1} = H^1((X/\Gamma_1)_s, \mathcal{F}) \arrow[r, hook] & H^1((X/\Gamma_1)_{\bar{\eta}}, \mathcal{F}) \arrow[d, equal] \\
& H^1(X_{\bar{\eta}}, \mathcal{F})^{\mathrm{SL} \cap \Gamma_1}
\end{tikzcd}
\]
induces an isomorphism of $\Sigma$-components
\[
H^1((X/\Gamma_1)_s, \mathcal{F})(\Sigma) \xrightarrow{\ \sim\ } H^1((X/\Gamma_1)_{\bar{\eta}}, \mathcal{F})(\Sigma) ;
\]
in particular the subspace
\[
H^1((X/\Gamma_1)_{\bar{\eta}}, \mathcal{F})(\Sigma)
\]
has good reduction over the field $(L)^{\det(\Gamma_1)}$.
\end{enumerate}
\end{corollary}

\begin{proof}
Simply apply the exact functor $M \mapsto M(\Sigma)$ to the commutative diagram
\[
\begin{tikzcd}[row sep=2.5em, column sep=2em]
H^1((X/\Gamma_1)_s, \mathcal{F}) \arrow[r, "\sim"] \arrow[d, hook] & H^1(X_s, \mathcal{F})^{\Gamma_1} \arrow[r, hook] \arrow[d, hook] & H^1(X_s, \mathcal{F}) \arrow[d, hook] \\
H^1((X/\Gamma_1)_{\bar{\eta}}, \mathcal{F}) \arrow[r, "\sim"] & H^1(X_{\bar{\eta}}, \mathcal{F})^{\mathrm{SL} \cap \Gamma_1} \arrow[r, hook] & H^1(X_{\bar{\eta}}, \mathcal{F}) . \qedhere
\end{tikzcd}
\]
\end{proof}

\section{Application to modular curves: explicitation of the numerical criterion}
\label{sec:14.3}

\kmitem\label{item:14.3.1}
Fix an integer $N \ge 1$ prime to $p$, and a subgroup
\[
\Gamma \subset \GL(2, \mathbb{Z}/N\mathbb{Z}) .
\]
Denote by $B$ the ring
\[
B = (\mathbb{Z}[\zeta_N])^{\det(\Gamma)} \subset \mathbb{Z}[\zeta_N] .
\]
We suppose that the moduli problem
\[
\mathcal{P} = ([\Gamma(N)]/\Gamma)^{B\text{-}\mathrm{can}} \tens{B} B[1/N] \quad\text{on } (\mathrm{Ell}/B[1/N])
\]
is representable. As we have seen, the compactified moduli scheme
\[
\begin{tikzcd}[row sep=1.8em]
\overline{\mathfrak{M}}(\mathcal{P}) \arrow[d] \\
\Spec(B[1/N])
\end{tikzcd}
\]
is a proper and smooth curve over $B[1/N]$ with geometrically connected fibers, in which
\[
(\overline{\mathfrak{M}}(\mathcal{P}) - \mathfrak{M}(\mathcal{P}))^{\mathrm{red}} = \mathrm{Cusps}(\mathcal{P})
\]
is finite etale over $B[1/N]$.

\kmitem\label{item:14.3.2}
Let us temporarily denote by
\begin{equation}\label{eq:14.3.2.1}
\begin{tikzcd}[row sep=1.8em, column sep=2.5em]
E_{\mathrm{univ}} \arrow[d, "f_{\mathrm{univ}}"] \\
\mathfrak{M}(\mathcal{P}) \arrow[r, "j", hook] & \overline{\mathfrak{M}}(\mathcal{P})
\end{tikzcd}
\end{equation}
the universal elliptic curve carried by $\mathfrak{M}(\mathcal{P})$. Fix a prime $\ell$ dividing $N$ (to assure that $\ell$ is invertible on $\overline{\mathfrak{M}}(\mathcal{P})$, and that $\ell \neq p$). Then
\begin{equation}\label{eq:14.3.2.2}
R^1 (f_{\mathrm{univ}})_* \mathbb{Q}_\ell
\end{equation}
is a lisse $\mathbb{Q}_\ell$ sheaf of rank two on $\mathfrak{M}(\mathcal{P})$, and on every geometric fiber of $\overline{\mathfrak{M}}(\mathcal{P}) \to \Spec(B[1/N])$, its geometric monodromy group is known to be an open subgroup of $\mathrm{SL}(2, \mathbb{Q}_\ell)$. This sheaf is automatically tamely ramified along the cusps.

\begin{proposition}\label{prop:14.3.3}
If $-1$ operates without fixed points on the space of cusp-labels for $\Gamma \cap \mathrm{SL}(2, \mathbb{Z}/N\mathbb{Z})$, then the geometric local monodromy at the cusps of $R^1(f_{\mathrm{univ}})_* \mathbb{Q}_\ell$ on $\mathfrak{M}(\mathcal{P})$ is unipotent of infinite order.
\end{proposition}

\begin{proof}
Because the sheaf $R^1(f_{\mathrm{univ}})_* \mathbb{Q}_\ell$ is tame along the cusps, and the cusps are finite etale over the base $\Spec(B[1/N])$ which is connected, it suffices to prove this on a single geometric fiber of $\overline{\mathfrak{M}}(\mathcal{P}) \to \Spec(B[1/N])$. Choose the complex embedding
\[
B[1/N] \mathrel{\lhook\joinrel\longrightarrow} \mathbb{Z}[\zeta_N, 1/N] \mathrel{\lhook\joinrel\longrightarrow} \mathbb{C}, \quad \zeta_N \mapsto \exp(2\pi i/N) ,
\]
and work on the corresponding $\overline{\mathfrak{M}}(\mathcal{P}) \tens{} \mathbb{C}$. Then we have
\[
\overline{\mathfrak{M}}(\mathcal{P}) \tens{B[1/N]} \mathbb{C} = \overline{\mathfrak{M}}(\mathcal{P} \tens{B[1/N]} \mathbb{Z}[\zeta_N, 1/N]) \tens{\mathbb{Z}[\zeta_N, 1/N]} \mathbb{C} ,
\]
so we may assume, to begin with, that $\Gamma \subset \mathrm{SL}(2, \mathbb{Z}/N\mathbb{Z})$.

On the complex manifold $(\overline{\mathfrak{M}}(\mathcal{P}) \tens{} \mathbb{C})^{\mathrm{an}}$, we have, canonically,
\[
(R^1(f_{\mathrm{univ}})_* \mathbb{Q}_\ell)^{\mathrm{an}} \simeq (R^1(f_{\mathrm{univ}}^{\mathrm{an}})_* \mathbb{Z}) \tens{} \mathbb{Q}_\ell ,
\]
so it suffices to look at the integral local monodromy at the cusps. At each cusp, the local monodromy is given by the transpose inverse of an element
\[
A \in \mathrm{SL}(2, \mathbb{Z}), \quad A \bmod N \text{ lies in } \Gamma .
\]
The element $A^2$ is unipotent of infinite order (because over a double covering of a punctured neighborhood of each cusp, the universal curve becomes a Tate curve $\mathrm{Tate}(q)$ viewed over $\mathbb{C}((q^{1/f}))$, for some divisor $f$ of $2N$).

Therefore $A$ has eigenvalues $\pm 1$, necessarily equal because $A$ lies in $\mathrm{SL}(2, \mathbb{Z})$. So
\[
A = \pm U
\]
for some unipotent $U \in \mathrm{SL}(2, \mathbb{Z})$. If $A = U$, we are done.

If not, in a suitable basis of $(\mathbb{Z})^2$, we have
\[
A = \begin{pmatrix} -1 & a \\ 0 & -1 \end{pmatrix} .
\]
In this case, the cusp-label $(1,0)$ satisfies
\[
(1,0)A = -(1,0) .
\]
Because $A \bmod N$ lies in $\Gamma$, this says that $-1$ fixes $(1,0)$ viewed as a cusp-label for $\Gamma$.
\end{proof}

\noindent\textbf{Remark.} Recall (\ref{cor:10.13.8}) that this same condition on $-1$ guarantees the existence of $\underline{\omega}$ on $\overline{\mathfrak{M}}(\mathcal{P})$. The two conditions, namely existence of $\underline{\omega}$ on $\overline{\mathfrak{M}}(\mathcal{P})$ and unipotence at infinity of $R^1(f_{\mathrm{univ}})_* \mathbb{Q}_\ell$, seem to be equivalent, but we do not know how to prove this.

\kmitem\label{item:14.3.4}
For the rest of this chapter, we will assume that $\mathcal{P}$ satisfies the additional hypothesis:
\begin{equation}\label{eq:14.3.4.1}
\pm 1 \text{ operates freely on the cusp-labels for } \Gamma \cap \mathrm{SL}(2, \mathbb{Z}/N\mathbb{Z}) .
\end{equation}
Then for every $d \ge 0$, the sheaf
\begin{equation}\label{eq:14.3.4.2}
\mathcal{F}(d) \overset{\mathrm{dfn}}{=} j_*(\mathrm{Symm}^d(R^1(f_{\mathrm{univ}})_* \mathbb{Q}_\ell)) \quad\text{on } \overline{\mathfrak{M}}(\mathcal{P})
\end{equation}
is lisse on $\mathfrak{M}(\mathcal{P})$ of rank $d+1$, and its restriction to $\mathrm{Cusps}(\mathcal{P})$ is a lisse sheaf of rank one. Because we are in mixed characteristic, $\mathcal{F}(d)$ is automatically tamely ramified along the cusps.

We will make constant use of the following \emph{fundamental property} of the sheaves $\mathcal{F}(d)$.

\begin{customprop}{14.3.4.3}\label{prop:14.3.4.3}
Let $k$ be an algebraically closed field,
\[
B[1/N] \to k
\]
a ring homomorphism, $C$ a proper smooth connected curve over $k$,
\[
\begin{tikzcd}[row sep=1.8em]
C \arrow[d, "\pi"] \\
\overline{\mathfrak{M}}(\mathcal{P}) \tens{} k
\end{tikzcd}
\]
a finite morphism, $D \subset C$ the inverse image of the cusps,
\[
C - D \mathrel{\lhook\joinrel\xrightarrow{\quad j\quad}} C
\]
the inclusion. Then
\begin{enumerate}[label=(\arabic*)]
\item For $d = 0$, $\mathcal{F}(0)$ is the constant sheaf $\mathbb{Q}_\ell$, $H^0(C, \mathbb{Q}_\ell) \simeq \mathbb{Q}_\ell$, $H^2(C, \mathbb{Q}_\ell) \simeq \mathbb{Q}_\ell(-1)$.
\item For $d > 0$, $\pi^*\mathcal{F}(d) \xrightarrow{\ \sim\ } j_* j^* \pi^*\mathcal{F}(d)$ is lisse on $C - D$, and tame along $D$. It is an absolutely irreducible representation of $\pi_1(C - D)$, and we have $H^0(C, \pi^*\mathcal{F}(d)) = H^2(C, \pi^*\mathcal{F}(d)) = 0$.
\end{enumerate}
\end{customprop}

\begin{proof}
This is immediate from the fact that $\mathcal{F}(1)$ on $\overline{\mathfrak{M}}(\mathcal{P}) \tens{} k$ has geometric monodromy open in $\mathrm{SL}(2, \mathbb{Q}_\ell)$, and unipotent non-trivial local monodromy around each cusp; because $\pi_1(C - D) \to \pi_1(\overline{\mathfrak{M}}(\mathcal{P}) \tens{} k)$ has image of finite index, the same statements hold for $\pi^*(\mathcal{F}(1))$ on $C$ itself. The symmetric power $\mathrm{Symm}^d$, $d \ge 1$, of any two-dimensional representation whose image is open in $\mathrm{SL}(2, \mathbb{Q}_\ell)$ is irreducible and non-trivial. The symmetric powers of any non-trivial unipotent in $\mathrm{SL}(2, \mathbb{Q}_\ell)$ all have a single Jordan block, so one-dimensional spaces of invariants.
\end{proof}

\kmitem\label{item:14.3.5}
We will now apply the theory of the preceding section to this situation. Fix a ring homomorphism
\[
B[1/N] \to k = \text{an algebraic closure of } \mathbb{F}_p .
\]
Because $B[1/N]$ is finite etale over $\mathbb{Z}[1/N]$, this extends uniquely to a ring homomorphism
\[
B[1/N] \to W(k) .
\]
We take
\begin{align*}
V &= W(k), \quad\text{fraction field denoted } K \\
Y &= \overline{\mathfrak{M}}(\mathcal{P}) \tens{B[1/N]} W(k) \\
\mathcal{F}(d) \text{ on } Y &= \text{the sheaf } \mathcal{F}(d) \text{ on } \overline{\mathfrak{M}}(\mathcal{P}), \text{ pulled back to } Y \\
L &= \text{the galois extension } K(\zeta_{p^n}) \text{ of } K, \text{ with} \\
&\quad\ \text{galois group } G = \Gal(L/K) \xrightarrow{\ \sim\ } (\mathbb{Z}/p^n\mathbb{Z})^\times
\end{align*}
\[
\begin{tikzcd}[row sep=1.8em]
X = \text{the scheme } \overline{\mathfrak{M}}(\mathcal{P}, [\mathrm{bal.}\ \Gamma_1(p^n)]^{\mathrm{can}}) \tens{B[1/N, \zeta_{p^n}]} (V_L) \arrow[d] \\
Y \tens{V} V_L
\end{tikzcd}
\]
\begin{align*}
\Gamma &= \text{the group } (\mathbb{Z}/p^n\mathbb{Z})^\times \times (\mathbb{Z}/p^n\mathbb{Z})^\times, \text{ operating} \\
&\quad\ \text{modularly on } X \text{ by its standard action} \\
&\quad\ (P, Q) \xmapsto{\alpha, \beta} (\alpha P, \beta Q) \\
&\quad\ \text{on the moduli problem } [\mathrm{bal.}\ \Gamma_1(p^n)] \\
\Gamma \xrightarrow{\det} G &\simeq (\mathbb{Z}/p^n\mathbb{Z})^\times \text{ is } (\alpha, \beta) \mapsto \alpha\beta .
\end{align*}
Then all the hypotheses of the general theorem (\ref{thm:14.2.4}) are satisfied. Because $\Gamma$ is \emph{abelian}, we may take for $\Gamma_2$ the anti-diagonal subgroup
\[
\Delta = \{(\alpha, \alpha^{-1}) \mid \alpha \in (\mathbb{Z}/p^n\mathbb{Z})^\times\} \subset \Gamma ,
\]
which operates $Y \tens{V} V_L$-linearly on $X$, and commutes with the $\det$-semilinear action of $\Gamma$.

The quotient
\[
X/(1 \times (\mathbb{Z}/p^n\mathbb{Z})^\times)
\]
is none other than the $Y$-scheme
\[
\begin{tikzcd}[row sep=1.8em]
\overline{\mathfrak{M}}(\mathcal{P}, [\Gamma_1(p^n)]) \tens{B[1/N]} W(k) \arrow[d] \\
Y = \overline{\mathfrak{M}}(\mathcal{P}) \tens{B[1/N]} W(k) ,
\end{tikzcd}
\]
and the action of $\Delta$ on this quotient is by the usual action
\[
P \xmapsto{(\alpha, \alpha^{-1})} \alpha P
\]
of $(\mathbb{Z}/p^n\mathbb{Z})^\times$ on the moduli problem $[\Gamma_1(p^n)]$.

The generic fiber of the quotient
\[
X/(1 \times (\mathbb{Z}/p^n\mathbb{Z})^\times)
\]
is just the usual descent of $X \tens{V_L} L$ to a smooth curve over $K$. For any subgroup
\[
H \subset (\mathbb{Z}/p^n\mathbb{Z})^\times ,
\]
the generic fiber of the quotient
\[
X/1 \times H
\]
is the usual descent of $X \tens{V_L} L$ to a smooth curve over $L^H$, namely
\begin{align*}
(X/(1 \times H)) \tens{(V_L)^H} L^H &= (X/(1 \times (\mathbb{Z}/p^n\mathbb{Z})^\times)) \tens{V} K \tens{K} (L^H) \\
&= \overline{\mathfrak{M}}(\mathcal{P}, [\Gamma_1(p^n)]) \tens{B[1/N]} L^H .
\end{align*}
Applying the numerical criterion to this situation, with $\Gamma_1 = 1 \times H$, we find

\begin{theorem}\label{thm:14.3.6}
Let $\Sigma$ be any $\mathbb{Q}$-rational set of $\overline{\mathbb{Q}}$-characters of the group $\Delta \simeq (\mathbb{Z}/p^n\mathbb{Z})^\times$, and let $H \subset (\mathbb{Z}/p^n\mathbb{Z})^\times$ be a subgroup. A sufficient condition that the $I_K$-stable subspace
\[
H^1(\overline{\mathfrak{M}}(\mathcal{P}, [\Gamma_1(p^n)]) \tens{B[1/N]} \overline{K}, \mathcal{F}(d))(\Sigma)
\]
have good reduction over the field $L^H = K(\zeta_{p^n})^H$ is that the inequality of dimensions
\begin{multline*}
\dim \Bigl( H^1(\overline{\mathfrak{M}}(\mathcal{P}, [\mathrm{bal.}\ \Gamma_1(p^n)]^{\mathrm{can}}) \tens{B[1/N, \zeta_{p^n}]} k, \mathcal{F}(d))^{1 \times H}(\Sigma) \Bigr) \\
\le \dim \Bigl( H^1(\overline{\mathfrak{M}}(\mathcal{P}, [\Gamma_1(p^n)]) \tens{B[1/N]} \overline{K}, \mathcal{F}(d))(\Sigma) \Bigr)
\end{multline*}
be an equality.
\end{theorem}

\begin{keyremark}\label{rem:14.3.7}
Hypotheses and notations as in (\ref{item:14.3.1}, \ref{item:14.3.5}, \ref{thm:14.3.6}) above, let us denote by $K_0$ the fraction field of $B$, and by $L_0$ the finite extension of $K_0$ defined by
\[
L_0 = (K_0(\zeta_{p^n}))^H .
\]
Fix an embedding $L_0 \hookrightarrow \overline{K}$ extending the chosen $K_0 \hookrightarrow K$, and denote by $\mathfrak{p}$ the corresponding $p$-adic place of $L_0$, and by
\[
D_\mathfrak{p} \subset \Gal(\overline{L}_0/L_0)
\]
the decomposition group at $\mathfrak{p}$.

Because the group $1 \times H$ acts $L_0$-linearly on the $B[1/N, \zeta_{p^n}]$-scheme $\overline{\mathfrak{M}}(\mathcal{P}, [\mathrm{bal.}\ \Gamma_1(p^n)]^{\mathrm{can}})$, while the scheme $\overline{\mathfrak{M}}(\mathcal{P}, [\Gamma_1(p^n)])$ itself lives over $B[1/N]$, it follows that the cohomology group
\[
H^1(\overline{\mathfrak{M}}(\mathcal{P}, [\Gamma_1(p^n)]) \tens{B[1/N]} \overline{K}, \mathcal{F}(d))
\]
is itself a representation of the entire galois group $\Gal(\overline{K}_0/K_0)$, and that the inclusion
\[
\begin{tikzcd}[row sep=1.8em]
H^1(\overline{\mathfrak{M}}(\mathcal{P}, [\mathrm{bal.}\ \Gamma_1(p^n)]^{\mathrm{can}}) \tens{B[1/N, \zeta_{p^n}]} k, \mathcal{F}(d))^{1 \times H} \arrow[d, hook] \\
H^1(\overline{\mathfrak{M}}(\mathcal{P}, [\Gamma_1(p^n)]) \tens{B[1/N]} \overline{K}, \mathcal{F}(d))
\end{tikzcd}
\]
is equivariant for the action of the subgroup
\[
D_\mathfrak{p} \subset \Gal(\overline{L}_0/L_0) \subset \Gal(\overline{K}_0/K_0) ,
\]
with the inertia subgroup $I_\mathfrak{p} \subset D_\mathfrak{p}$ operating trivially on the source.

Therefore, if this inclusion induces an isomorphism of $\Sigma$-components, the resulting isomorphism of $\Sigma$-components is an isomorphism in the category of unramified representations of $D_\mathfrak{p}$.
\end{keyremark}

\section{Characters and conductors}
\label{sec:14.4}

\kmitem\label{item:14.4.1} Fix an integer $n \ge 1$, and a subgroup
\[
H_0 \subset (\mathbb{Z}/p^n\mathbb{Z})^\times
\]
which satisfies the two following conditions:
\begin{enumerate}[label=\arabic*)]
\item $H_0$ has order dividing $p-1$ (so $H_0 = \{1\}$ if $p = 2$)
\item if $n = 1$, $H_0$ is not the entire group $(\mathbb{Z}/p\mathbb{Z})^\times$.
\end{enumerate}

\kmitem\label{item:14.4.2} A $\overline{\mathbb{Q}}$-valued character
\[
\chi : \Delta \simeq (\mathbb{Z}/p^n\mathbb{Z})^\times \longrightarrow (\overline{\mathbb{Q}})^\times
\]
is said to be of \emph{type $(H_0, v)$} for $v$ an integer $1 \le v \le n$ if it satisfies the following three conditions:
\begin{enumerate}[label=\arabic*)]
\item $\chi(H_0) = 1$
\item $\chi(x) = 1$ if $x \equiv 1 \pmod{p^v}$
\item there exists an element $y \in (\mathbb{Z}/p^n\mathbb{Z})^\times$ with $y \equiv 1 \pmod{p^{v-1}}$, and $\chi(y) \ne 1$.
\end{enumerate}
For example, the characters of type $(H_0, n)$ are exactly the characters of maximal conductor which kill $H_0$.

Every non-trivial $\overline{\mathbb{Q}}$-valued $\chi$ which is trivial on $H_0$ is of type $(H_0, v)$ for the unique integer $v \in [1,n]$ such that $\chi$ has conductor $p^v$. When convenient, we will say that the \emph{trivial} character is of type $(H_0, 0)$.

\kmitem\label{item:14.4.3} For each $v \in [1,n]$, the set of $\overline{\mathbb{Q}}$-valued characters of type $(H_0, v)$ is $\mathbb{Q}$-rational. In any linear representation of the group $\Delta \simeq (\mathbb{Z}/p^n\mathbb{Z})^\times$, on a vector space $M$ over a field of characteristic zero, we denote by
\begin{equation}\label{eq:14.4.3.1}
M(H_0, v)
\end{equation}
the sum of the $\chi$-components for all $\chi$ of type $(H_0, v)$. Clearly we have an isomorphism of $\Delta$-representations
\begin{equation}\label{eq:14.4.3.2}
M(H_0, v) \simeq \begin{cases}
(M)^{H_0 \times (1+(p^v))} / (M)^{H_0 \times (1+(p^{v-1}))} & \text{if } v \ge 2 \\[10pt]
(M)^{H_0 \times (1+(p))} / (M)^\Delta & \text{if } v = 1 .
\end{cases}
\end{equation}

\section{The Good Reduction Theorem}
\label{sec:14.5}

\begin{theorem}\label{thm:14.5.1}
Let $N \ge 1$ be an integer prime to $p$, $\Gamma \subset \GL(2, \mathbb{Z}/N\mathbb{Z})$ a subgroup, $B$ the ring $(\mathbb{Z}[\zeta_N])^{\det(\Gamma)}$. Suppose that the moduli problem
\[
\mathcal{P} = ([\Gamma(N)]/\Gamma) \tens{\mathbb{Z}[\zeta_N]} B[1/N] \quad \text{on } (\Ell/B[1/N])
\]
is representable, and that $-1$ operates without fixed points on the cusp-labels of $\Gamma \cap \SL(2, \mathbb{Z}/N\mathbb{Z})$ (cf.~(\ref{lem:10.7.1})).

Let $\ell$ be a prime number dividing $N$, $d$ an integer $\ge 0$, and $\mathcal{F}(d)$ the $\mathbb{Q}_\ell$-sheaf on $\overline{\mathfrak{M}}(\mathcal{P})$ constructed above (\ref{eq:14.3.4.2}).

Let $n \ge 1$, $H_0 \subset (\mathbb{Z}/p^n\mathbb{Z})^\times$ a proper subgroup of order dividing $p-1$, and $1 \le v \le n$ an integer. Suppose that we are in one of the following two situations:
\begin{enumerate}[label=\alph*)]
\item $v = n$
\item $v > [n/2]$, and $H_0$ is trivial.
\end{enumerate}
For any ring homomorphism
\[
B[1/N] \longrightarrow W(k), \quad k \text{ an algebraic closure of } \mathbb{F}_p ,
\]
denoting by $K$ the fraction field of $W(k)$, the subspace
\[
H^1(\overline{\mathfrak{M}}(\mathcal{P}, [\Gamma_1(p^n)]) \tens{B[1/N]} K, \mathcal{F}(d))(H_0, v)
\]
has good reduction over the field
\[
K(\zeta_{p^n})^{H_0 \times (1+(p^v))} .
\]
\end{theorem}

\begin{proof}
We will prove the theorem by verifying the numerical criterion of good reduction.

We introduce some notation:
\begin{equation}\label{eq:14.5.1.1}
\begin{cases}
X_{\bar{\eta}} = \overline{\mathfrak{M}}(\mathcal{P}, [\Gamma_1(p^n)]) \tens{B[1/N]} \overline{K} \\[6pt]
Y_{\bar{\eta}} = \overline{\mathfrak{M}}(\mathcal{P}) \tens{B[1/N]} \overline{K} \\[6pt]
X_s = \overline{\mathfrak{M}}(\mathcal{P}, [\mathrm{bal.}\ \Gamma_1(p^n)]^{\mathrm{can}}) \tens{B[1/N, \zeta_{p^n}]} k \\[6pt]
Y_s = \overline{\mathfrak{M}}(\mathcal{P}) \tens{B[1/N]} k
\end{cases}
\end{equation}
and for each pair of non-negative integers $(a,b)$ with $a+b=n$,
\begin{equation}\label{eq:14.5.1.2}
X_s(a,b) = \overline{\mathfrak{M}}\bigl( (\mathcal{P} \tens{B[1/N]} k)^{(\sigma^{-b})}, [\mathrm{Ig}(p^{\max(a,b)})] \bigr) ,
\end{equation}
viewed as a $Y_s$-scheme by $\mathrm{pr}_b$ (cf.~(\ref{item:12.10.5})).

For each integer $v$ in $[0,n]$, we denote by
\[
H(v) \subset (\mathbb{Z}/p^n\mathbb{Z})^\times
\]
the subgroup
\begin{equation}\label{eq:14.5.1.3}
H(v) = \begin{cases}
(\mathbb{Z}/p^n\mathbb{Z})^\times & \text{if } v = 0 \\[4pt]
H_0 \times (1+(p^v)) & \text{if } 1 \le v \le n .
\end{cases}
\end{equation}
We denote by
\[
\Delta(H(v)) \subset \Delta
\]
the corresponding subgroup of $\Delta \simeq (\mathbb{Z}/p^n\mathbb{Z})^\times$:
\begin{equation}\label{eq:14.5.1.4}
\Delta(H(v)) = \{ (\alpha, \alpha^{-1}) \mid \alpha \in H(v) \} \subset \Delta .
\end{equation}

We will prove that the inclusion coming from the vanishing cycle exact sequence is in fact an equality
\refstepcounter{kmitem}%
\begin{equation}\label{eq:14.5.2}\tag{\thekmitem}
H^1(X_s, \mathcal{F}(d))^{1 \times H(v)}(H_0, v) = H^1(X_{\bar{\eta}}, \mathcal{F}(d))(H_0, v) .
\end{equation}

\kmitem\label{item:14.5.3} We first explain why it is sufficient to verify this equality under the additional hypothesis that $\Gamma$ lies in $\SL(2, \mathbb{Z}/N\mathbb{Z})$.

The equality in question depends only on the inverse image of $\mathcal{F}(d)$ on the $W(k)[\zeta_{p^n}]$-scheme
\[
\overline{\mathfrak{M}}(\mathcal{P} \tens{B[1/N]} W(k), [\mathrm{bal.}\ \Gamma_1(p^n)]^{\mathrm{can}}) \tens{\mathbb{Z}[\zeta_{p^n}]} W(k)[\zeta_{p^n}] .
\]
Because $\mathbb{Z}[\zeta_N, 1/N]$ is finite etale over $B[1/N]$, there exists an extension of the given homomorphism $B[1/N] \to W(k)$
\[
\begin{tikzcd}[row sep=2.5em, column sep=2.5em]
\mathbb{Z}[\zeta_N, 1/N] \arrow[r] & W(k) \\
B[1/N] \arrow[u, hook] \arrow[ur] &
\end{tikzcd}
\]
For any such extension, we have
\begin{align*}
\mathcal{P} \tens{B[1/N]} W(k) &= (\mathcal{P} \tens{B[1/N]} \mathbb{Z}[\zeta_N, 1/N]) \tens{\mathbb{Z}[\zeta_N, 1/N]} W(k) \\
&= \bigl( ([\Gamma(N)]/\Gamma \cap \SL(2, \mathbb{Z}/N\mathbb{Z}))^{\mathbb{Z}[\zeta_N]\text{-}\mathrm{can}} \bigr) \tens{\mathbb{Z}[\zeta_N]} W(k) .
\end{align*}
Thus we may replace $\Gamma$ by $\Gamma \cap \SL(2, \mathbb{Z}/N\mathbb{Z})$, and $B[1/N]$ by $\mathbb{Z}[\zeta_N, 1/N]$, without changing what is to be proven. In this case, we have by hypothesis that $-1$ has no fixed points on the cusp-labels of $\Gamma \cap \SL(2, \mathbb{Z}/N\mathbb{Z}) = \Gamma$. This permits us to make use of (\ref{thm:13.12.3}). Recall also that by (\ref{cor:10.13.8}), the line bundle $\underline{\omega}$ exists on $\overline{\mathfrak{M}}(\mathcal{P})$.

We will verify---eventually---the numerical criterion in this case.

\kmitem\label{item:14.5.4} To verify the numerical criterion, we must show that
\begin{equation}\label{eq:14.5.4.1}
\dim \Bigl( H^1(X_s, \mathcal{F}(d))^{1 \times H(v)}(H_0, v) \Bigr) = \dim \Bigl( H^1(X_{\bar{\eta}}, \mathcal{F}(d))(H_0, v) \Bigr) .
\end{equation}
In view of the explicit description (\ref{eq:14.4.3.2}) of the $(H_0, v)$ component of a representation of $\Delta$, this amounts to proving that
\begin{multline}\label{eq:14.5.4.2}
\dim \left( \frac{H^1(X_s, \mathcal{F}(d))^{\Delta(H(v)) \times (1 \times H(v))}}{H^1(X_s, \mathcal{F}(d))^{\Delta(H(v-1)) \times (1 \times H(v))}} \right) \\
= \dim \left( \frac{H^1(X_{\bar{\eta}}, \mathcal{F}(d))^{\Delta(H(v))}}{H^1(X_{\bar{\eta}}, \mathcal{F}(d))^{\Delta(H(v-1))}} \right) .
\end{multline}

\kmitem\label{item:14.5.5} We next explain how to reduce this to an equality involving Euler characteristics. By the exact sequence of vanishing cycles
\begin{multline*}
0 \longrightarrow H^1(X_s, \mathcal{F}(d)) \longrightarrow H^1(X_{\bar{\eta}}, \mathcal{F}(d)) \\
\longrightarrow \bigoplus_{x\ \mathrm{s.s.}} (\mathcal{F}(d)_x \tens{\mathbb{Q}_\ell} (\mathbb{Q}_\ell^I/\mathbb{Q}_\ell)) \longrightarrow 0
\end{multline*}
the group $1 \times H(v)$ operates \emph{trivially} on the vanishing cycles (they are supported at the supersingular points, each of which is fixed by $(\mathbb{Z}/p^n\mathbb{Z})^\times \times (\mathbb{Z}/p^n\mathbb{Z})^\times$). Therefore the quotient space
\begin{equation}\label{eq:14.5.5.1}
\frac{H^1(X_{\bar{\eta}}, \mathcal{F}(d))}{H^1(X_s, \mathcal{F}(d))}
\end{equation}
has $1 \times H(v)$ acting \emph{trivially}, and this action is via the diagonal embedding $\Delta \subset (\mathbb{Z}/p^n\mathbb{Z})^\times \times (\mathbb{Z}/p^n\mathbb{Z})^\times$. Therefore for any subgroup $H \subset (\mathbb{Z}/p^n\mathbb{Z})^\times$, the invariants under $\Delta(H) \times (1 \times H(v))$ in the vanishing cycles are the same as the invariants under $\Delta(H)$. Thus taking invariants in the short exact sequence of representations of $\Delta$
\begin{equation}\label{eq:14.5.5.2}
0 \to H^1(X_s, \mathcal{F}(d)) \to H^1(X_{\bar{\eta}}, \mathcal{F}(d)) \to (\text{vanishing cycles}) \to 0
\end{equation}
gives an exact sequence of invariants under $\Delta(H(v)) \times (1 \times H(v))$, and similarly for $\Delta(H(v-1)) \times (1 \times H(v))$. Therefore the difference of dimensions in (\ref{eq:14.5.4.2}) is equal to the difference of the Euler characteristics:
\begin{multline}\label{eq:14.5.5.3}
\chi(X_s / (\Delta(H(v)) \times (1 \times H(v))), \mathcal{F}(d)) \\
- \chi(X_s / (\Delta(H(v-1)) \times (1 \times H(v))), \mathcal{F}(d)) \\
= \chi(X_{\bar{\eta}} / \Delta(H(v)), \mathcal{F}(d)) - \chi(X_{\bar{\eta}} / \Delta(H(v-1)), \mathcal{F}(d)) .
\end{multline}
To see this, we check that the $H^0$ and $H^2$ terms do not contribute to the difference.

If $d > 0$, then because $\mathcal{F}(d) = \mathrm{Sym}^d(\mathcal{F}(1))$ and $\mathcal{F}(1)$ is the relative $H^1$ of the universal elliptic curve, $\mathcal{F}(d)$ has no non-zero global sections or $H^2$ on any irreducible component of $X_{\bar{\eta}}$ or $X_s$, so for $d > 0$ we have
\begin{equation}\label{eq:14.5.5.4}
\begin{cases}
H^0(X_{\bar{\eta}}, \mathcal{F}(d)) = H^2(X_{\bar{\eta}}, \mathcal{F}(d)) = 0 \\[4pt]
H^0(X_s, \mathcal{F}(d)) = H^2(X_s, \mathcal{F}(d)) = 0 .
\end{cases}
\end{equation}

If $d = 0$, $\mathcal{F}(0) = \mathbb{Q}_\ell$. Then $H^0(X_{\bar{\eta}}, \mathbb{Q}_\ell) = \mathbb{Q}_\ell$, with trivial action of $\Delta$, so it contributes identically to both terms on the right. Similarly, $H^2(X_{\bar{\eta}}, \mathbb{Q}_\ell) \simeq \mathbb{Q}_\ell(-1)$ with trivial action of $\Delta$, so it also contributes identically to both terms on the right.

As for $X_s$, its connected components are also its irreducible components, and the action of $(\mathbb{Z}/p^n\mathbb{Z})^\times \times (\mathbb{Z}/p^n\mathbb{Z})^\times$ on $H^0(X_s, \mathbb{Q}_\ell)$ and on $H^2(X_s, \mathbb{Q}_\ell)$ factors through the permutation of the components. We will show that both
\begin{equation}\label{eq:14.5.5.5}
\begin{cases}
\Delta(H(v)) \times (1 \times H(v)) \\[4pt]
\Delta(H(v-1)) \times (1 \times H(v))
\end{cases}
\end{equation}
have the \emph{same} orbits on the components of $X_s$. Indeed, the components are parameterized by pairs $(a,b)$ with $a+b=n$ and an element of $(\mathbb{Z}/p^{\min(a,b)}\mathbb{Z})^\times$. The group $(\mathbb{Z}/p^n\mathbb{Z})^\times \times (\mathbb{Z}/p^n\mathbb{Z})^\times$ acts by
\begin{equation}\label{eq:14.5.5.6}
(u_1, u_2) : (a,b, \zeta) \mapsto (a,b, (u_1/u_2) \zeta \bmod p^{\min(a,b)}) .
\end{equation}
If $n = 1$, the components are $(1,0)$ and $(0,1)$, each occurring with multiplicity one ($\min(a,b) = 0$), so the group acts trivially on the set of components!

If $n \ge 2$, we are assuming that either $v = n \ge 2$ or $v > [n/2] \ge 1$, so in all cases
\[
v-1 \ge 1 .
\]
In this case, the groups
\[
\begin{cases}
\Delta(H(v)) \times (1 \times H(v)) \\[4pt]
\Delta(H(v-1)) \times (1 \times H(v))
\end{cases}
\]
both contain the subgroup $1 \times H_0$, and both are contained in the group $H(v-1) \times H(v-1) = (H_0 \times (1+(p^{v-1}))) \times (H_0 \times (1+(p^{v-1})))$. So both operate on the space of components through the image of $1 \times H_0$, i.e.,
\begin{equation}\label{eq:14.5.5.7}
\begin{aligned}
H^2(X_s, \mathbb{Q}_\ell)^{1 \times H_0} &= H^2(X_s, \mathbb{Q}_\ell)^{\Delta(H(v)) \times (1 \times H(v))} \\
&= H^2(X_s, \mathbb{Q}_\ell)^{\Delta(H(v-1)) \times (1 \times H(v))} .
\end{aligned}
\end{equation}

So again in the case $d = 0$, the numerical criterion boils down to the equality of Euler characteristics (\ref{eq:14.5.5.3}), this time with $d = 0$.

\kmitem\label{item:14.5.6} We now divide each Euler characteristic into two terms, one from the cusps and the second from the complement of the cusps. Thus let us denote by
\[
\begin{cases}
CX_{\bar{\eta}} = \text{the cusps of } X_{\bar{\eta}} \\[4pt]
CY_{\bar{\eta}} = \text{the cusps of } Y_{\bar{\eta}} \\[4pt]
CX_s = \text{the cusps of } X_s \\[4pt]
CY_s = \text{the cusps of } Y_s .
\end{cases}
\]
The decompositions
\begin{align*}
X_{\bar{\eta}} &= CX_{\bar{\eta}} \cup (X_{\bar{\eta}} - CX_{\bar{\eta}}) \\
X_s &= CX_s \cup (X_s - CX_s)
\end{align*}
are respected by the groups of operators $\Delta$, $(\mathbb{Z}/p^n\mathbb{Z})^\times \times (\mathbb{Z}/p^n\mathbb{Z})^\times$ respectively (because these groups operate as $Y_{\bar{\eta}}$ and $Y_s$-automorphisms respectively, and these decompositions of $X_{\bar{\eta}}$, $X_s$ are the inverse images of the analogous decompositions of $Y_{\bar{\eta}}$ and $Y_s$ respectively).

\begin{proposition}\label{prop:14.5.7}
Under the hypotheses of the Main Theorem (\ref{thm:14.5.1}), the ``cuspidal part'' of the numerical criterion holds; we have, for any $d \ge 0$,
\begin{multline*}
\chi(CX_s / (\Delta(H(v)) \times (1 \times H(v))), \mathcal{F}(d)) \\
- \chi(CX_s / (\Delta(H(v-1)) \times (1 \times H(v))), \mathcal{F}(d)) \\
= \chi(CX_{\bar{\eta}} / \Delta(H(v)), \mathcal{F}(d)) - \chi(CX_{\bar{\eta}} / \Delta(H(v-1)), \mathcal{F}(d)) .
\end{multline*}
\end{proposition}

\begin{proof}
The restriction of $\mathcal{F}(d)$ to the cusps is lisse of rank one, so this amounts to a count of the cusps: we need
\begin{multline}\label{eq:14.5.7.1}
\#(CX_s / (\Delta(H(v)) \times (1 \times H(v)))) \\
- \#(CX_s / (\Delta(H(v-1)) \times (1 \times H(v)))) \\
= \#(CX_{\bar{\eta}} / \Delta(H(v))) - \#(CX_{\bar{\eta}} / \Delta(H(v-1))) .
\end{multline}
We know that the scheme of cusps in
\[
\overline{\mathfrak{M}}(\mathcal{P}, [\mathrm{bal.}\ \Gamma(p^n)]^{\mathrm{can}})
\]
is finite etale over $\mathbb{Z}[1/N, \zeta_N, \zeta_{p^n}]$, a disjoint union of sections. Therefore we have a canonical bijection of finite sets
\[
CX_{\bar{\eta}} \xrightarrow{\ \sim\ } CX_s
\]
which is equivariant for the action of $\Delta$.

Therefore we may ``forget about'' $X_{\bar{\eta}}$; we must prove the equality
\begin{multline}\label{eq:14.5.7.2}
\#(CX_s / (\Delta(H(v)) \times (1 \times H(v)))) \\
- \#(CX_s / (\Delta(H(v-1)) \times (1 \times H(v)))) \\
= \#(CX_s / \Delta(H(v))) - \#(CX_s / \Delta(H(v-1))) .
\end{multline}
But the cusps of $X_s$ are the disjoint union of the cusps on the various irreducible components (in the notation of (\ref{eq:14.5.1.2}))
\begin{equation}\label{eq:14.5.7.3}
\coprod_{a+b=n} (\mathbb{Z}/p^{\min(a,b)}\mathbb{Z})^\times \text{ copies of } X_s(a,b) .
\end{equation}
For any subgroup (cf.~(\ref{item:13.12.2}))
\[
K \subset (\mathbb{Z}/p^n\mathbb{Z})^\times \times (\mathbb{Z}/p^n\mathbb{Z})^\times ,
\]
recall that we defined subgroups
\begin{equation}\label{eq:14.5.7.4}
K \supset K(a,b) = \{ (k_1, k_2) \in K \text{ with } k_1 \equiv k_2 \bmod p^{\min(a,b)} \}
\end{equation}
\begin{multline}\label{eq:14.5.7.5}
(\mathbb{Z}/p^{\min(a,b)}\mathbb{Z})^\times \supset \mathrm{Comp}\ K(a,b) \\
= \{ k_2/k_1 \bmod p^{\min(a,b)}, \text{ for } (k_1, k_2) \in K \}
\end{multline}
\begin{equation}\label{eq:14.5.7.6}
\begin{aligned}
&(\mathbb{Z}/p^{\max(a,b)}\mathbb{Z})^\times \supset \mathrm{Im}\ K(a,b) \\
&= \begin{cases}
\begin{aligned}
&\{ k_1 \bmod p^{\max(a,b)} \text{ such that } \\
&\quad (k_1, k_2) \in K(a,b) \text{ for some } k_2 \} \quad \text{if } a \le b
\end{aligned} \\[10pt]
\begin{aligned}
&\{ k_2 \bmod p^{\max(a,b)} \text{ such that } \\
&\quad (k_1, k_2) \in K(a,b) \text{ for some } k_1 \} \quad \text{if } a \ge b
\end{aligned}
\end{cases}
\end{aligned}
\end{equation}
and we proved (\ref{thm:13.12.3}) that the irreducible components of $X_s/K$ outside the supersingular points are
\begin{equation}\label{eq:14.5.7.7}
\coprod_{a+b=n} \left( \begin{array}{@{}c@{}} (\mathbb{Z}/p^{\min(a,b)}\mathbb{Z})^\times / \mathrm{Comp}\ K(a,b) \text{ copies of} \\[4pt] X_s(a,b) / \mathrm{Im}\ K(a,b) \end{array} \right) .
\end{equation}
We have seen (\ref{thm:12.7.1}) that the entire group $(\mathbb{Z}/p^{\max(a,b)}\mathbb{Z})^\times$ operates freely on $CX_s(a,b)$, and the cusps of $X_s(a,b)$ are in fact a $(\mathbb{Z}/p^{\max(a,b)}\mathbb{Z})^\times$ torseur over the cusps of $Y_s^{(\sigma^{-b})}$. So for any subgroup
\[
K \subset (\mathbb{Z}/p^n\mathbb{Z})^\times \times (\mathbb{Z}/p^n\mathbb{Z})^\times ,
\]
we have the following formula:
\begin{equation}\label{eq:14.5.7.8}
\#(CX_s / K) = \sum_{a+b=n} \frac{\phi(p^{\min(a,b)}) \phi(p^{\max(a,b)})}{\#(\mathrm{Comp}\ K(a,b)) \cdot \#(\mathrm{Im}\ K(a,b))} \#(CY_s) .
\end{equation}
We will apply this formula, one $(a,b)$-term at a time, to the four different $K$'s we are juggling. The computation is straightforward, and we give the result in the form of tables: in the second table, $H_0^{(2)}$ denotes the subgroup of squares of elements of $H_0$.

\par\addvspace{6pt plus 2pt minus 1pt}\noindent\phantomsection\label{item:14.5.7.9}(14.5.7.9)\enspace \emph{Table, if $n = 1$; $K \to$}:
\begin{center}
\resizebox{\linewidth}{!}{%
\begin{tabular}{|l|c|c|c|c|}
\hline
& $\Delta(H(1)) = \Delta(H_0)$ & $\Delta(H(0)) = \Delta$ & $\begin{array}{@{}c@{}}\Delta(H(1)) \times (1 \times H(1))\\= H_0 \times H_0\end{array}$ & $\Delta(H(0)) \times (1 \times H_0)$ \\
\hline
$\mathrm{Comp}\ K(0,1)$ & $1$ & $1$ & $1$ & $1$ \\
\hline
$\mathrm{Im}\ K(0,1)$ & $H_0$ & $\mathbb{F}_p^\times$ & $H_0$ & $\mathbb{F}_p^\times$ \\
\hline
$\mathrm{Comp}\ K(1,0)$ & $1$ & $1$ & $1$ & $1$ \\
\hline
$\mathrm{Im}\ K(1,0)$ & $H_0$ & $\mathbb{F}_p^\times$ & $H_0$ & $\mathbb{F}_p^\times$ \\
\hline
\end{tabular}%
}
\end{center}

\par\addvspace{6pt plus 2pt minus 1pt}\noindent\phantomsection\label{item:14.5.7.10}(14.5.7.10)\enspace \emph{Table, if $n \ge 2$; $K \to$ (so $v-1 \ge 1$)}:
\begin{center}
\resizebox{\linewidth}{!}{%
\begin{tabular}{|l|c|c|c|c|}
\hline
& $\Delta(H(v))$ & $\Delta(H(v-1))$ & $\Delta(H(v)) \times (1 \times H(v))$ & $\Delta(H(v-1)) \times (1 \times H(v))$ \\
\hline
\parbox[c]{2.5cm}{$\mathrm{Comp}\ K(a,b)$\\ if $a \ge b \ge 1$} & $H_0^{(2)}$ & $H_0^{(2)}$ & $H_0$ & $H_0$ \\
\hline
\parbox[c]{2.5cm}{$\mathrm{Im}\ K(a,b)$\\ if $a \ge b \ge 1$} & $\begin{array}{@{}c@{}}(H_0 \cap \pm 1)\\\times (1+(p^v))\\\bmod p^a\end{array}$ & $\begin{array}{@{}c@{}}(H_0 \cap \pm 1)\\\times (1+(p^{v-1}))\\\bmod p^a\end{array}$ & $H(v) \bmod p^a$ & $\begin{array}{@{}c@{}}H_0 \times (1+(p^{v-1}))\\= H(v-1) \bmod p^a\end{array}$ \\
\hline
\parbox[c]{2.5cm}{$\mathrm{Comp}\ K(a,b)$\\ if $b \ge a \ge 1$} & $H_0^{(2)}$ & $H_0^{(2)}$ & $H_0$ & $H_0$ \\
\hline
\parbox[c]{2.5cm}{$\mathrm{Im}\ K(a,b)$\\ $b \ge a \ge 1$} & $\begin{array}{@{}c@{}}(H_0 \cap \pm 1)\\\times (1+(p^v))\\\bmod p^b\end{array}$ & $\begin{array}{@{}c@{}}(H_0 \cap \pm 1)\\\times (1+(p^{v-1}))\\\bmod p^b\end{array}$ & $H(v) \bmod p^b$ & $\begin{array}{@{}c@{}}H_0 \times (1+(p^{v-1}))\\= H(v-1) \bmod p^b\end{array}$ \\
\hline
$\mathrm{Comp}\ K(0,n)$ & $1$ & $1$ & $1$ & $1$ \\
\hline
$\mathrm{Im}\ K(0,n)$ & $H(v)$ & $H(v-1)$ & $H(v)$ & $H(v-1)$ \\
\hline
$\mathrm{Comp}\ K(n,0)$ & $1$ & $1$ & $1$ & $1$ \\
\hline
$\mathrm{Im}\ K(n,0)$ & $H(v)$ & $H(v-1)$ & $H(v)$ & $H(v-1)$ \\
\hline
\end{tabular}%
}
\end{center}

In this second table, $n \ge 2$, we now consider separately the two cases
\[
\begin{cases}
v = n, \quad H_0 \text{ an arbitrary subgroup of order dividing } p-1 \\
v > [n/2], \quad H_0 \text{ trivial} .
\end{cases}
\]
In the first of these cases, $v = n$, so $v-1 \ge n-1 \ge \max(a,b)$ whenever both $a,b$ are $\ge 1$. Then $H(n) = H_0$, and the table boils down to

\par\addvspace{6pt plus 2pt minus 1pt}\noindent\phantomsection\label{item:14.5.7.11}(14.5.7.11)\enspace \emph{Table, if $n \ge 2$ and $v = n$}:
\begin{center}
\resizebox{\linewidth}{!}{%
\begin{tabular}{|l|c|c|c|c|}
\hline
& $\Delta(H_0)$ & $\Delta(H(n-1))$ & $\Delta(H_0) \times (1 \times H_0)$ & $\Delta(H(n-1)) \times (1 \times H_0)$ \\
\hline
\parbox[c]{2.5cm}{$\mathrm{Comp}\ K(a,b)$\\ if $a, b \text{ both } \ge 1$} & $H_0^{(2)}$ & $H_0^{(2)}$ & $H_0$ & $H_0$ \\
\hline
\parbox[c]{2.5cm}{$\mathrm{Im}\ K(a,b)$\\ if $a, b \text{ both } \ge 1$} & $H_0 \cap (\pm 1)$ & $H_0 \cap (\pm 1)$ & $H_0$ & $H_0$ \\
\hline
\parbox[c]{2.5cm}{$\mathrm{Comp}\ K(a,b)$\\ if $a \text{ or } b = 0$} & $1$ & $1$ & $1$ & $1$ \\
\hline
\parbox[c]{2.5cm}{$\mathrm{Im}\ K(a,b)$\\ if $a \text{ or } b = 0$} & $H_0$ & $H(n-1)$ & $H_0$ & $H(n-1)$ \\
\hline
\end{tabular}%
}
\end{center}

In the second of these cases, $v > [n/2] \ge 1$, $H_0$ trivial, the table becomes even simpler:

\par\addvspace{6pt plus 2pt minus 1pt}\noindent\phantomsection\label{item:14.5.7.12}(14.5.7.12)\enspace \emph{Table, if $n \ge 2$, $v > [n/2]$, and $H_0$ trivial}:
\begin{center}
\resizebox{\linewidth}{!}{%
\begin{tabular}{|l|c|c|c|c|}
\hline
& $\Delta(H(v))$ & $\Delta(H(v-1))$ & $\Delta(H(v)) \times (1 \times H(v))$ & $\Delta(H(v-1)) \times (1 \times H(v))$ \\
\hline
$\mathrm{Comp}\ K(a,b)$ & $1$ & $1$ & $1$ & $1$ \\
\hline
$\mathrm{Im}\ K(a,b)$ & $H(v) \bmod p^{\max(a,b)}$ & $H(v-1) \bmod p^{\max(a,b)}$ & $H(v) \bmod p^{\max(a,b)}$ & $H(v-1) \bmod p^{\max(a,b)}$ \\
\hline
\end{tabular}%
}
\end{center}

Now let us return to the formula (\ref{eq:14.5.7.8})
\[
\#(CX_s / K) = \sum_{a+b=n} \frac{\phi(p^{\min(a,b)}) \phi(p^{\max(a,b)})}{\#(\mathrm{Comp}\ K(a,b)) \cdot \#(\mathrm{Im}\ K(a,b))} \#(CY_s) ,
\]
and temporarily christen the individual terms
\setcounter{equation}{12}%
\begin{equation}\label{eq:14.5.7.13}
\#(CX_s / K)(a,b) \overset{\mathrm{dfn}}{=} \frac{\phi(p^{\min(a,b)}) \phi(p^{\max(a,b)})}{\#(\mathrm{Comp}\ K(a,b)) \cdot \#(\mathrm{Im}\ K(a,b))} \#(CY_s) .
\end{equation}
We claim that for \emph{each} $(a,b)$ with $a+b=n$, we have an equality of the individual $(a,b)$-components of the formula we are to prove, namely we claim
\begin{multline}\label{eq:14.5.7.14}
\#(CX_s / \Delta(H(v)))(a,b) - \#(CX_s / \Delta(H(v-1)))(a,b) \\
= \#(CX_s / (\Delta(H(v)) \times (1 \times H(v))))(a,b) \\
- \#(CX_s / (\Delta(H(v-1)) \times (1 \times H(v))))(a,b) .
\end{multline}
Indeed, a look at the tables shows that in the case
\[
n \ge 2, \quad v = n, \quad H_0 \text{ arbitrary}, \quad a, b \text{ both } \ge 1
\]
both sides of the alleged equality vanish, and it shows that in all the other cases we must consider, both sides of the alleged equality are \emph{term-by-term} equal.
\end{proof}

\kmitem\label{item:14.5.8} We now return to the problem of proving the Main Theorem by proving the equality of Euler characteristics for every $\mathcal{F}(d)$, $d \ge 0$ (cf.~(\ref{eq:14.5.5.3})):
\begin{multline}\label{eq:14.5.8.1}
\chi(X_s / (\Delta(H(v)) \times (1 \times H(v))), \mathcal{F}(d)) \\
- \chi(X_s / (\Delta(H(v-1)) \times (1 \times H(v))), \mathcal{F}(d)) \\
= \chi(X_{\bar{\eta}} / \Delta(H(v)), \mathcal{F}(d)) - \chi(X_{\bar{\eta}} / \Delta(H(v-1)), \mathcal{F}(d)) .
\end{multline}
The sheaf $\mathcal{F}(d)$ is lisse of rank $d+1$ outside the cusps, tamely ramified along the cusps, and its restriction to the cusps is lisse of rank one. Therefore each of the four Euler characteristics we are considering, written generically as $\chi(Z, \mathcal{F}(d))$, satisfies
\begin{equation}\label{eq:14.5.8.2}
\begin{aligned}
\chi(Z, \mathcal{F}(d)) &= \chi(Z - CZ, \mathcal{F}(d)) + \chi(CZ, \mathcal{F}(d)) \\
&= (d+1) \chi(Z - CZ) + \#(CZ) .
\end{aligned}
\end{equation}
Since we have already verified the cuspidal part of the equality, we are reduced to proving the following equality of topological Euler characteristics:
\begin{multline}\label{eq:14.5.8.3}
\chi((X_s - CX_s) / (\Delta(H(v)) \times (1 \times H(v)))) \\
- \chi((X_s - CX_s) / (\Delta(H(v-1)) \times (1 \times H(v)))) \\
= \chi((X_{\bar{\eta}} - CX_{\bar{\eta}}) / \Delta(H(v))) - \chi((X_{\bar{\eta}} - CX_{\bar{\eta}}) / \Delta(H(v-1))) .
\end{multline}

\kmitem\label{item:14.5.9} \emph{The right-hand difference}

The entire group $\Delta \simeq (\mathbb{Z}/p^n\mathbb{Z})^\times$ operates freely on $X_{\bar{\eta}} - CX_{\bar{\eta}}$ (e.g., outside of $p$, $[\Gamma_1(p^n)]$ is finite etale over $[\Gamma_0(p^n)]$ with group $(\mathbb{Z}/p^n\mathbb{Z})^\times$), and because we are in characteristic zero we have, for any subgroup $G \subset (\mathbb{Z}/p^n\mathbb{Z})^\times \simeq \Delta$,
\begin{equation}\label{eq:14.5.9.1}
\chi((X_{\bar{\eta}} - CX_{\bar{\eta}})/G) = \left( \frac{1}{\# G} \right) \chi(X_{\bar{\eta}} - CX_{\bar{\eta}})
\end{equation}
so the right-hand difference of Euler characteristics is equal to
\begin{equation}\label{eq:14.5.9.2}
\left( \frac{1}{\# H(v)} - \frac{1}{\# H(v-1)} \right) \chi(X_{\bar{\eta}} - CX_{\bar{\eta}}) .
\end{equation}
Now the covering
\begin{equation}\label{eq:14.5.9.3}
\begin{tikzcd}[row sep=2.5em]
X_{\bar{\eta}} - CX_{\bar{\eta}} \arrow[d] \\
Y_{\bar{\eta}} - CY_{\bar{\eta}}
\end{tikzcd}
\end{equation}
is itself finite etale of degree $p^{2n} - p^{2n-2}$ (the degree of $[\Gamma_1(p^n)]$ over $(\Ell)$), so we may rewrite the right-hand difference of Euler characteristics as
\begin{equation}\label{eq:14.5.9.4}
p^{2n-2}(p^2-1) \left( \frac{1}{\# H(v)} - \frac{1}{\# H(v-1)} \right) \chi(Y_{\bar{\eta}} - CY_{\bar{\eta}}) .
\end{equation}

\kmitem\label{item:14.5.10} \emph{The left-hand difference}

We next analyze the left-hand difference of Euler characteristics. The curve $X_s - CX_s$ is fully ramified over $Y_s - CY_s$ at each supersingular point, so the supersingular points of $X_s - CX_s$ map isomorphically to the supersingular points of $Y_s - CY_s$. Therefore the supersingular points of $X_s - CX_s$ are individually fixed by the entire group $(\mathbb{Z}/p^n\mathbb{Z})^\times \times (\mathbb{Z}/p^n\mathbb{Z})^\times$, simply because this group operates by $Y_s$-automorphisms. Therefore in the left-hand difference of Euler characteristics, the contribution of the supersingular points to both terms is the same, namely $\#\mathrm{SS}(Y_s)$, so we may rewrite the left-hand difference as
\begin{multline}\label{eq:14.5.10.1}
\chi((X_s - CX_s - \mathrm{SS}(X_s)) / (\Delta(H(v)) \times (1 \times H(v)))) \\
- \chi((X_s - CX_s - \mathrm{SS}(X_s)) / (\Delta(H(v-1)) \times (1 \times H(v)))) .
\end{multline}
The open curve
\begin{equation}\label{eq:14.5.10.2}
X_s - CX_s - \mathrm{SS}(X_s)
\end{equation}
is the \emph{disjoint} union of irreducible curves, namely
\begin{equation}\label{eq:14.5.10.3}
\coprod_{a+b=n} (\mathbb{Z}/p^{\min(a,b)}\mathbb{Z})^\times \text{ copies of } X_s(a,b) - CX_s(a,b) - \mathrm{SS} X_s(a,b) .
\end{equation}
For any subgroup
\begin{equation}\label{eq:14.5.10.4}
K \subset (\mathbb{Z}/p^n\mathbb{Z})^\times \times (\mathbb{Z}/p^n\mathbb{Z})^\times ,
\end{equation}
the quotient
\begin{equation}\label{eq:14.5.10.5}
(X_s - CX_s - \mathrm{SS}(X_s))/K
\end{equation}
is the disjoint union
\begin{equation}\label{eq:14.5.10.6}
\coprod_{a+b=n} \left( \begin{array}{@{}c@{}} \displaystyle \frac{\phi(p^{\min(a,b)})}{\#\mathrm{Comp}\ K(a,b)} \text{ copies of} \\[10pt] (X_s(a,b) - CX_s(a,b) - \mathrm{SS} X_s(a,b)) / \mathrm{Im}\ K(a,b) \end{array} \right) .
\end{equation}
We now consider in detail the two groups
\begin{equation}\label{eq:14.5.10.7}
\begin{cases}
K_1 = \Delta(H(v)) \times (1 \times H(v)) \\[4pt]
K_2 = \Delta(H(v-1)) \times (1 \times H(v)) .
\end{cases}
\end{equation}
For ease of reference, we repeat the salient parts of our earlier tables, in the relevant cases:
\begin{center}
\resizebox{0.9\linewidth}{!}{%
\begin{tabular}{l|c|c}
& $K_1$ & $K_2$ \\
\hline
\underline{if $n = 1$} & & \\
$\mathrm{Comp}\ K(a,b)$ & $1$ & $1$ \\
\hline
$\mathrm{Im}\ K(a,b)$ & $H_0$ & $\mathbb{F}_p^\times$ \\
\hline\hline
\underline{if $v = n \ge 1$} & & \\
$\mathrm{Comp}\ K(a,b)$ if $a, b \text{ both } \ge 1$ & $H_0$ & $H_0$ \\
\hline
$\mathrm{Im}\ K(a,b)$ if both $a, b \ge 1$ & $H_0$ & $H_0$ \\
\hline
$\mathrm{Comp}\ K(a,b)$ if $a \text{ or } b = 0$ & $1$ & $1$ \\
\hline
$\mathrm{Im}\ K(a,b)$ if $a \text{ or } b = 0$ & $H_0$ & $H(n-1)$ \\
\hline\hline
\underline{if $v > [n/2] \ge 1$, $H_0$ trivial} & & \\
$\mathrm{Comp}\ K(a,b)$ & $1$ & $1$ \\
\hline
$\mathrm{Im}\ K(a,b)$ & $H(v) \bmod p^{\max(a,b)}$ & $H(v-1) \bmod p^{\max(a,b)}$ \\
\hline
\end{tabular}%
}
\end{center}
We will consider separately the three cases.

\kmitem\label{item:14.5.11} \emph{Case~I ($n = 1$)}

In this case
\begin{align*}
&\chi((X_s - CX_s - \mathrm{SS}(X_s))/K_1) - \chi(\text{''}/K_2) \\
&= \sum_{a+b=1} \Bigl[ \chi((X_s(a,b) - CX_s(a,b) - \mathrm{SS} X_s(a,b))/H_0) - \chi(\text{''}/\mathbb{F}_p^\times) \Bigr] .
\end{align*}
Because $H_0$ and $\mathbb{F}_p^\times$ have order prime to $p$, and operate freely on the Igusa curve minus its supersingular points, we may rewrite this
\begin{align*}
&\sum_{a+b=1} \left( \frac{1}{\# H_0} - \frac{1}{\#\mathbb{F}_p^\times} \right) \chi(X_s(a,b) - CX_s(a,b) - \mathrm{SS} X_s(a,b)) \\
&= 2 \left( \frac{1}{\# H_0} - \frac{1}{\#\mathbb{F}_p^\times} \right) \chi(X_s(1,0) - CX_s(1,0) - \mathrm{SS} X_s(1,0)) .
\end{align*}
Because $X_s(1,0)$ is the Igusa curve of level $p$, it minus its supersingular points is an etale $\mathbb{F}_p^\times$-torsor over $Y_s - \mathrm{SS}(Y_s)$, so we may rewrite the last expression as
\[
2(p-1) \left( \frac{1}{\# H_0} - \frac{1}{\#\mathbb{F}_p^\times} \right) \chi(Y_s - CY_s - \mathrm{SS} Y_s) .
\]

\kmitem\label{item:14.5.12} \emph{Case~II ($v = n \ge 2$)}

In this case all the $(a,b)$ terms with both $a$ and $b \ge 1$ are identical for $K_1$ and $K_2$; the only terms which contribute to the difference are $(a,b) = (n,0)$ or $(0,n)$. For each of these terms, the \emph{difference} is
\[
\chi((X_s(n,0) - CX_s(n,0) - \mathrm{SS} X_s(n,0))/H_0) - \chi(\text{''}/H(n-1)) .
\]
The group $H_0$ operates freely on the open Igusa curve, and has order prime to $p$. The group $H(n-1)$ is the product of $H_0$ and the group of order $p$, $1+(p^{n-1})$; it also operates freely. So we may ``isolate'' the $H_0$-contribution, and rewrite the above difference as
\[
\frac{1}{\# H_0} \Bigl[ \chi(X_s(n,0) - CX_s(n,0) - \mathrm{SS} X_s(n,0)) - \chi(\text{''} / (1+(p^{n-1}))) \Bigr] .
\]
Since both sides have the \emph{same} number of supersingular points, namely $\#(\mathrm{SS} Y_s)$, we may rewrite this difference as
\begin{multline*}
\frac{1}{\# H_0} \Bigl[ \chi(X_s(n,0) - CX_s(n,0)) \\
- \chi((X_s(n,0) - CX_s(n,0))/(1+(p^{n-1}))) \Bigr] .
\end{multline*}
But the quotient of $X_s(n,0)$ by $1+(p^a)$ is $X_s(a,0)$, so we may rewrite this as
\[
\frac{1}{\# H_0} \Bigl[ \chi(X_s(n,0) - CX_s(n,0)) - \chi(X_s(n-1,0) - CX_s(n-1,0)) \Bigr] .
\]
But this expression is the difference provided by both the $(n,0)$ term and the $(0,n)$ term, so all together the left-hand difference of Euler characteristics is
\[
\frac{2}{\# H_0} \Bigl[ \chi(X_s(n,0) - CX_s(n,0)) - \chi(X_s(n-1,0) - CX_s(n-1,0)) \Bigr] .
\]

\par\addvspace{6pt plus 2pt minus 1pt}\noindent\phantomsection\label{item:14.5.12b}(14.5.12)\enspace\emph{Case~III ($v > [n/2] \ge 1$, $H_0$ trivial)}

In this case, the left-hand difference is
\begin{align*}
&\sum_{a+b=n} \phi(p^{\min(a,b)}) \Bigl[ \chi((X_s(a,b) - CX_s(a,b) - \mathrm{SS}(X_s(a,b)))/H(v)) \\
&\qquad\qquad\qquad\quad - \chi(\text{''} / H(v-1)) \Bigr] .
\end{align*}
The groups $H(v), H(v-1)$ operate through their images $\bmod p^{\max(a,b)}$. If $v-1 \ge \max(a,b)$, then both groups become \emph{trivial} $\bmod p^{\max(a,b)}$, so the only terms which contribute a non-zero difference are those with
\[
\max(a,b) \ge v .
\]
For each such term, the curve $X_s(a,b)$ is an Igusa curve of level $p^{\max(a,b)}$, whose quotient by $H(v) = 1+(p^v)$ (resp.~by $H(v-1) = 1+(p^{v-1})$) is the Igusa curve of level $p^v$ (resp.~$p^{v-1}$). Moreover, all these curves have the same number of supersingular points, namely $\#\mathrm{SS} Y_s$. All in all, the left-hand difference is
\begin{multline*}
\sum_{\substack{a+b=n \\ \max(a,b) \ge v}} \phi(p^{\min(a,b)}) \Bigl[ \chi(X_s(v,0) - CX_s(v,0)) \\
- \chi(X_s(v-1,0) - CX_s(v-1,0)) \Bigr] .
\end{multline*}

\kmitem\label{item:14.5.13} \emph{Recapitulation of what remains to prove}

For ease of reference, we put together the formulas we are reduced to verifying:

\par\addvspace{6pt plus 2pt minus 1pt}\noindent\phantomsection\label{item:14.5.13.1}(14.5.13.1)\enspace\emph{Case~I ($n = 1$)}:
\begin{multline*}
(p^2-1) \left( \frac{1}{\# H_0} - \frac{1}{\#\mathbb{F}_p^\times} \right) \chi(Y_{\bar{\eta}} - CY_{\bar{\eta}}) \\
= 2(p-1) \left( \frac{1}{\# H_0} - \frac{1}{\#\mathbb{F}_p^\times} \right) \chi(Y_s - CY_s - \mathrm{SS} Y_s) .
\end{multline*}

\par\addvspace{6pt plus 2pt minus 1pt}\noindent\phantomsection\label{item:14.5.13.2}(14.5.13.2)\enspace\emph{Case~II ($v = n \ge 2$)}:
\begin{multline*}
\left( \frac{p^{2n-2}(p^2-1)}{\# H_0} \right) \left( 1 - \frac{1}{p} \right) \chi(Y_{\bar{\eta}} - CY_{\bar{\eta}}) \\
= \frac{2}{\# H_0} \Bigl[ \chi(X_s(n,0) - CX_s(n,0)) \\
- \chi(X_s(n-1,0) - CX_s(n-1,0)) \Bigr] .
\end{multline*}

\par\addvspace{6pt plus 2pt minus 1pt}\noindent\phantomsection\label{item:14.5.13.3}(14.5.13.3)\enspace\emph{Case~III ($v > [n/2] \ge 1$)}:
\begin{multline*}
p^{2n-2}(p^2-1) \left( \frac{1}{p^{n-v}} - \frac{1}{p^{n+1-v}} \right) \chi(Y_{\bar{\eta}} - CY_{\bar{\eta}}) \\
= \sum_{\substack{a+b=n \\ \max(a,b) \ge v}} \phi(p^{\min(a,b)}) \Bigl[ \chi(X_s(v,0) - CX_s(v,0)) \\
- \chi(X_s(v-1,0) - CX_s(v-1,0)) \Bigr] .
\end{multline*}

\kmitem\label{item:14.5.14} \emph{End of the proof: Case~I}

We must show that
\[
(p+1) \chi(Y_{\bar{\eta}} - CY_{\bar{\eta}}) = 2 \chi(Y_s - CY_s - \mathrm{SS} Y_s) .
\]
Now $Y = \overline{\mathfrak{M}}(\mathcal{P}) \tens{\mathbb{Z}} W(k)$ has good reduction, and its cusps are finite etale over the base, so
\[
\chi(Y_{\bar{\eta}} - CY_{\bar{\eta}}) = \chi(Y_s - CY_s) ,
\]
and we are reduced to proving
\[
(p-1) \chi(Y_s - CY_s) = -2\#\mathrm{SS} Y_s .
\]
In terms of the line bundle $\underline{\omega}$ on $Y_s$ (which does indeed exist, by hypothesis), we have
\[
\chi(Y_s - CY_s) = -2 \deg(\underline{\omega})
\]
by (\ref{cor:10.13.12}), and we have
\[
\#\mathrm{SS} Y_s = (p-1) \deg(\underline{\omega})
\]
by Igusa's theorem (cf.~(\ref{thm:12.4.3})), whence the assertion. Q.E.D.

\kmitem\label{item:14.5.15} \emph{End of the proof: Case~II ($v = n \ge 2$)}

The formula to be proven in this case is
\begin{multline*}
p^{2n-3}(p^2-1)(p-1) \chi(Y_{\bar{\eta}} - CY_{\bar{\eta}}) \\
= 2 \Bigl[ \chi(X_s(n,0) - CX_s(n,0)) - \chi(X_s(n-1,0) - CX_s(n-1,0)) \Bigr] .
\end{multline*}
In terms of the line bundle $\underline{\omega}$ on $Y_s$, we have (cf.~(\ref{cor:10.13.12}))
\[
\chi(Y_{\bar{\eta}} - CY_{\bar{\eta}}) = -2 \deg(\underline{\omega}) ,
\]
and, for every $n \ge 1$ we have (cf.~(\ref{cor:12.9.4}))
\[
\chi(X_s(n,0) - CX_s(n,0)) = -p^n \phi(p^n) \deg(\underline{\omega}) .
\]
The formula to be shown boils down to the assertion that for $n \ge 2$,
\[
2 p^{2n-3}(p^2-1)(p-1) = 2 \Bigl( p^n \phi(p^n) - p^{n-1} \phi(p^{n-1}) \Bigr)
\]
whose verification we leave to the reader! Q.E.D.

\kmitem\label{item:14.5.16} \emph{End of the proof: Case~III ($v > [n/2] \ge 1$)}

The formula to be proven in this case is
\begin{multline*}
p^{n+v-3}(p^2-1)(p-1) \chi(Y_{\bar{\eta}} - CY_{\bar{\eta}}) \\
= \left( \sum_{\substack{a+b=n \\ \max(a,b) \ge v}} \phi(p^{\min(a,b)}) \right) \Bigl[ \chi(X_s(v,0) - CX_s(v,0)) \\
- \chi(X_s(v-1,0) - CX_s(v-1,0)) \Bigr] .
\end{multline*}
In the sum over $(a,b)$, the fact that $v > [n/2]$ means that the terms with $\max(a,b) \ge v$ occur in pairs,
\[
(v+r, n-v-r), \quad (n-v-r, v+r)
\]
for $r = 0, \dots, n-v$. So we have
\begin{align*}
\sum_{\substack{a+b=n \\ \max(a,b) \ge v}} \phi(p^{\min(a,b)}) &= 2 \sum_{r=0}^{n-v} \phi(p^{n-v-r}) \\
&= 2 \sum_{r=0}^{n-v} \phi(p^r) \\
&= 2 \#(\mathbb{Z}/p^{n-v}\mathbb{Z}) \\
&= 2 p^{n-v} ,
\end{align*}
and the formula to be proven becomes
\begin{multline*}
p^{2v-3}(p^2-1)(p-1) \chi(Y_{\bar{\eta}} - CY_{\bar{\eta}}) \\
= 2 \Bigl[ \chi(X_s(v,0) - CX_s(v,0)) - \chi(X_s(v-1,0) - CX_s(v-1,0)) \Bigr]
\end{multline*}
for $v \ge 2$. But this is exactly the formula we just checked in Case~II.
\end{proof}

\section{Explicitation of the Good Reduction Theorem}
\label{sec:14.6}

\kmitem\label{item:14.6.1} Throughout this section, we adopt the hypotheses and notations of the good reduction theorem (\ref{thm:14.5.1}).

With the notations $X_s, X_s(a,b), \dots$ of (\ref{eq:14.5.1.1}), we have proven that the inclusion of the vanishing cycle exact sequence
\[
H^1(X_s, \mathcal{F}(d)) \mathrel{\lhook\joinrel\longrightarrow} H^1(X_{\bar{\eta}}, \mathcal{F}(d))
\]
induces an equality
\refstepcounter{kmitem}%
\begin{equation}\label{eq:14.6.2}\tag{\thekmitem}
H^1(X_s, \mathcal{F}(d))^{1 \times H(v)}(H_0, v) = H^1(X_{\bar{\eta}}, \mathcal{F}(d))(H_0, v) .
\end{equation}

We wish to explicitly describe the cohomology group
\[
H^1(X_s, \mathcal{F}(d))^{1 \times H(v)}(H_0, v)
\]
in terms of the cohomology of certain Igusa curves.

Recall from (\ref{eq:14.5.5.2}) the cohomology exact sequence
\refstepcounter{kmitem}%
\begin{equation}\label{eq:14.6.3}\tag{\thekmitem}
\resizebox{\linewidth}{!}{%
\begin{tikzpicture}[>=to, node distance=0.8cm and 0.8cm]
\node (A) {$0$};
\node (B) [right=of A] {$H^0(X_s, \mathcal{F}(d))$};
\node (C) [right=of B] {$\displaystyle\bigoplus_{\substack{a+b=n \\ u \in (\mathbb{Z}/p^{\min(a,b)}\mathbb{Z})^\times}} H^0(X_s(a,b), \mathcal{F}(d))$};

\node (D) [below=1.5cm of A.west, anchor=west] {$\displaystyle\bigoplus_{\substack{\text{s.s. points} \\ y \text{ of } Y_s}} \mathcal{F}(d)_y \tens{\mathbb{Q}_\ell} (\mathbb{Q}_\ell^I/\mathbb{Q}_\ell)$};
\node (E) [right=of D] {$H^1(X_s, \mathcal{F}(d))$};
\node (F) [right=of E] {$\displaystyle\bigoplus_{\substack{a+b=n \\ u \in (\mathbb{Z}/p^{\min(a,b)}\mathbb{Z})^\times}} H^1(X_s(a,b), \mathcal{F}(d))$};
\node (G) [right=of F] {$0$.};

\draw[->] (A) -- (B);
\draw[->] (B) -- (C);
\draw[->] (D) -- (E);
\draw[->] (E) -- (F);
\draw[->] (F) -- (G);

\draw (C.east) to[out=0, in=0, looseness=2] ++(0, -0.6) -- ([yshift=0.6cm]D.west);
\draw[->] ([yshift=0.6cm]D.west) to[out=180, in=180, looseness=2] (D.west);
\end{tikzpicture}%
}
\end{equation}
Because $\mathcal{F}(d)$ is the inverse image of a sheaf on $Y_s$, while the entire group $(\mathbb{Z}/p^n\mathbb{Z})^\times \times (\mathbb{Z}/p^n\mathbb{Z})^\times$ operates $Y_s$-linearly on $X_s$, this entire group acts on the term
\[
\bigoplus_{y \in \mathrm{SS} Y_s} \mathcal{F}(d)_y \tens{\mathbb{Q}_\ell} (\mathbb{Q}_\ell^I/\mathbb{Q}_\ell)
\]
as $\bigoplus (\mathrm{id} \otimes (\text{permutation action on } I))$.

Therefore its $(H_0, v)$-component under the action of $\Delta$ is trivial (proof: by hypothesis $v-1 \ge [n/2] \ge \min(a,b)$ for any $a+b=n$, so the entire subgroup $(1+(p^{v-1}))^\times \times (1+(p^{v-1}))^\times$ of $(\mathbb{Z}/p^n\mathbb{Z})^\times \times (\mathbb{Z}/p^n\mathbb{Z})^\times$, and in particular its intersection with $\Delta$, acts trivially on the indexing set $I$ of irreducible components, because
\[
I = \coprod_{a+b=n} (\mathbb{Z}/p^{\min(a,b)}\mathbb{Z})^\times
\]
with action of $(\alpha, \beta)$ multiplication by $\beta/\alpha$). So we have an isomorphism
\refstepcounter{kmitem}%
\begin{align}\label{eq:14.6.4}\tag{\thekmitem}
&H^1(X_s, \mathcal{F}(d))(H_0, v) \xrightarrow{\ \sim\ } \nonumber \\
&\quad \left( \bigoplus_{\substack{a+b=n \\ u \in (\mathbb{Z}/p^{\min(a,b)}\mathbb{Z})^\times}} H^1(X_s(a,b), \mathcal{F}(d)) \right)(H_0, v) ,
\end{align}
and passing to $1 \times H(v)$ invariants as well we obtain an isomorphism
\refstepcounter{kmitem}%
\begin{align}\label{eq:14.6.5}\tag{\thekmitem}
&H^1(X_s, \mathcal{F}(d))^{1 \times H(v)}(H_0, v) \xrightarrow{\ \sim\ } \nonumber \\
&\quad \left( \bigoplus_{\substack{a+b=n \\ u \in (\mathbb{Z}/p^{\min(a,b)}\mathbb{Z})^\times}} H^1(X_s(a,b), \mathcal{F}(d)) \right)^{1 \times H(v)}(H_0, v) .
\end{align}
The entire group $(\mathbb{Z}/p^n\mathbb{Z})^\times \times (\mathbb{Z}/p^n\mathbb{Z})^\times$ respects the $(a,b)$-decomposition, so we may rewrite this
\refstepcounter{kmitem}%
\begin{equation}\label{eq:14.6.6}\tag{\thekmitem}
\resizebox{\linewidth}{!}{%
$\displaystyle
\begin{aligned}
&\bigoplus_{a+b=n} [(\mathbb{Z}/p^{\min(a,b)}\mathbb{Z})^\times \text{ copies of } H^1(X_s(a,b), \mathcal{F}(d))]^{1 \times H(v)}(H_0, v) \\
&\simeq \bigoplus_{a+b=n} \frac{[(\mathbb{Z}/p^{\min(a,b)}\mathbb{Z})^\times \text{ copies of } H^1(X_s(a,b), \mathcal{F}(d))]^{\Delta(H(v)) \times (1 \times H(v))}}{[\text{ \quad '' \quad }]^{\Delta(H(v-1)) \times (1 \times H(v))}} .
\end{aligned}
$%
}
\end{equation}
A glance at the table following (\ref{eq:14.5.10.7}) shows that only the terms $(a,b)$ with $\max(a,b) \ge v$ contribute, and that their contribution is
\refstepcounter{kmitem}%
\begin{multline}\label{eq:14.6.7}\tag{\thekmitem}
\bigoplus_{\substack{a+b=n \\ \max(a,b) \ge v}} (\mathbb{Z}/p^{\min(a,b)}\mathbb{Z})^\times \text{ copies of} \\
H^1(X_s(a,b)/(1+(p^v)), \mathcal{F}(d))(H_0, v) .
\end{multline}
The quotient
\refstepcounter{kmitem}%
\begin{equation}\label{eq:14.6.8}\tag{\thekmitem}
X_s(a,b)/(1+(p^v))
\end{equation}
is given explicitly as
\refstepcounter{kmitem}%
\begin{equation}\label{eq:14.6.9}\tag{\thekmitem}
\begin{aligned}
&\overline{\mathfrak{M}}((\mathcal{P} \tens{} k)^{(\sigma^{-b})}, [\mathrm{Ig}(p^{\max(a,b)})])/(1+(p^v)) \\
&= \overline{\mathfrak{M}}((\mathcal{P} \tens{} k)^{(\sigma^{-b})}, [\mathrm{Ig}(p^v)]) ,
\end{aligned}
\end{equation}
viewed as $\overline{\mathfrak{M}}(\mathcal{P} \tens{} k)$-scheme by $\mathrm{pr}_b$ (cf.~(\ref{item:12.10.3})).

We summarize these explicit results in the three cases.

\kmitem\label{item:14.6.10} \emph{Case~I, II ($v = n \ge 1$)}
\begin{multline*}
H^1(X_{\bar{\eta}}, \mathcal{F}(d))(H_0, n) \text{ has good reduction over } (K(\zeta_{p^n}))^{H_0} , \\
\begin{aligned}
&\simeq H^1(\overline{\mathfrak{M}}(\mathcal{P} \tens{} k, [\mathrm{Ig}(p^n)]), \mathcal{F}(d))(H_0, n) \\
&\quad \oplus H^1(\overline{\mathfrak{M}}((\mathcal{P} \tens{} k)^{(\sigma^{-n})}, [\mathrm{Ig}(p^n)]), \mathcal{F}(d))(H_0, n) .
\end{aligned}
\end{multline*}

\kmitem\label{item:14.6.11} \emph{Case~III ($v > [n/2] \ge 1, H_0$ trivial)}
\begin{multline*}
H^1(X_{\bar{\eta}}, \mathcal{F}(d))(H_0, v) \text{ has good reduction over } K(\zeta_{p^v}) \\
\simeq \bigoplus_{\substack{a+b=n, \max(a,b) \ge v \\ u \in (\mathbb{Z}/p^{\min(a,b)}\mathbb{Z})^\times}} H^1(\overline{\mathfrak{M}}((\mathcal{P} \tens{} k)^{(\sigma^{-b})}, [\mathrm{Ig}(p^v)]), \mathcal{F}(d))(H_0, v) .
\end{multline*}

\section{Application to Jacobians}
\label{sec:14.7}

\kmitem\label{item:14.7.1} Let $p$ be a prime number, $n \ge 1$, and $H_0 \subset (\mathbb{Z}/p^n\mathbb{Z})^\times$ a proper subgroup, of order dividing $p-1$. Let $N \ge 5$ be an integer prime to $p$ (so that the moduli problem $\mathcal{P} = [\Gamma_1(N)] \tens{} \mathbb{Z}[1/N]$ satisfies all the hypotheses of the good reduction theorem).

Let us denote by
\[
J_1(Np^n)
\]
the abelian variety over $\mathbb{Q}$ which is the \emph{Jacobian} of
\[
\overline{\mathfrak{M}}([\Gamma_1(Np^n)] \tens{} \mathbb{Q}) .
\]
This $J_1(Np^n)$ visibly has good reduction over $\mathbb{Z}[1/Np]$---simply because $\overline{\mathfrak{M}}([\Gamma_1(Np^n)] \tens{\mathbb{Z}} \mathbb{Z}[1/Np])$ is a proper smooth curve over $\mathbb{Z}[1/Np]$ with geometrically connected fibers, whose $\mathrm{Pic}^0$ is an abelian scheme over $\mathbb{Z}[1/Np]$ extending $J_1(Np^n)$.

The group $(\mathbb{Z}/p^n\mathbb{Z})^\times$ operates on $J_1(Np^n)$ through its modular action $P \mapsto aP$ on the moduli problem $[\Gamma_1(p^n)]$, so in the category of $\mathbb{Q}$-abelian varieties up to isogeny we can speak of the direct factor of type $(H_0, v)$ for any $1 \le v \le n$:
\[
J_1(Np^n)(H_0, v) \mathrel{\lhook\joinrel\longrightarrow} J_1(Np^n) .
\]

\begin{theorem}\label{thm:14.7.2}
\begin{enumerate}[label=(\arabic*)]
\item If $v = n$, the $\mathbb{Q}$-abelian variety up to isogeny $J_1(Np^n)(H_0, n)$ acquires good reduction at the unique $p$-adic place of the field $(\mathbb{Q}(\zeta_{p^n}))^{H_0}$, and its special fiber at this place is $\mathbb{F}_p$-isogenous to the direct sum of two copies of the $\mathbb{F}_p$-abelian variety up to isogeny
\[
\mathrm{Jac}\bigl( \overline{\mathfrak{M}}([\Gamma_1(N)] \tens{} \mathbb{F}_p, [\mathrm{Ig}(p^n)]) \bigr)(H_0, n) .
\]
\item If $v > [n/2] \ge 1$, and $H_0$ is trivial, the $\mathbb{Q}$-abelian variety up to isogeny $J_1(Np^n)(H_0, v)$ acquires good reduction at the unique $p$-adic place of $\mathbb{Q}(\zeta_{p^v})$, and its special fiber at this place is $\mathbb{F}_p$-isogenous to the direct sum
\[
\bigoplus_{\substack{a+b=n, \max(a,b) \ge v \\ u \in (\mathbb{Z}/p^{\min(a,b)}\mathbb{Z})^\times}} \mathrm{Jac}\bigl( \overline{\mathfrak{M}}([\Gamma_1(N)] \tens{} \mathbb{F}_p, [\mathrm{Ig}(p^v)]) \bigr)(H_0, v)
\]
($2 p^{n-v}$ copies of the isogeny-direct factor of $\mathrm{Jac}(\overline{\mathfrak{M}}([\Gamma_1(N)] \tens{} \mathbb{F}_p, [\mathrm{Ig}(p^v)]))$ of maximal conductor $p^v$ under the action of $(\mathbb{Z}/p^n\mathbb{Z})^\times$).
\end{enumerate}
\end{theorem}

\begin{proof}
This follows from the special case $d = 0$, $\mathcal{F}(d) = \text{the constant sheaf } \mathbb{Q}_\ell$, of the good reduction theorem, by the Neron-Ogg-Shafarevic criterion of good reduction (cf.~\cite{SeTa2}), the Key Remark (\ref{rem:14.3.7}), and Tate's isogeny theorem \cite{Ta1}.
\end{proof}
